% concatenated from Kollar_conjecture-ENS.tex

      \documentclass[11pt,a4paper,english]{smfart}

 \usepackage[margin=1.15in]{geometry}
%\usepackage{amscd,amssymb, amsmath, setspace, wasysym}
\usepackage{amscd,amssymb, amsmath,amsfonts, wasysym, mathrsfs, enumitem, stmaryrd, mathtools,hhline,color,enumitem, tikz-cd}
\usepackage{graphicx}
\usepackage[all, cmtip]{xy}
%\linespread{1.1} 
\usepackage[T1]{fontenc}
\usepackage{newtxtext}
\usepackage{newtxmath}
\usepackage{textcomp}
\usepackage{url}
\usepackage{array}

\usepackage{thmtools}
\usepackage{enumitem}
\usepackage[pagebackref=true]{hyperref} %linkcolor=violet,  citecolor=hot, urlcolor=magenta]
%make all the references and links clickable
\usepackage{nameref}
\usepackage{cleveref}
\usepackage{comment}
\usetikzlibrary{arrows,positioning,calc,intersections}
\setlist[itemize]{wide = 0pt, labelwidth = 2em, labelsep*=0em, itemindent = 0pt, leftmargin = \dimexpr\labelwidth + \labelsep\relax, noitemsep,topsep = 1ex,}
\setlist[enumerate]{wide = 0pt, labelwidth = 2em, labelsep*=0em, itemindent = 0pt, leftmargin = \dimexpr\labelwidth + \labelsep\relax, noitemsep,topsep = 1ex}


\definecolor{hot}{RGB}{65,105,225}


\renewcommand{\thefootnote}{(\textit{\alph{footnote}})}
%\UseRawInputEncoding
%\usepackage{refcheck}

% to make my stuff easier to write, i use \begin{thm} \end{thm} and here i define the shortcuts:
\theoremstyle{plain}
\newtheorem{theorem}{Theorem}[section]
\newtheorem{proposition}[theorem]{Proposition}
\newtheorem{thm_defn}[theorem]{Theorem/Definition} 
\newtheorem{thmx}{Theorem}
\renewcommand{\thethmx}{\Alph{thmx}}  
\newtheorem{corx}[thmx]{Corollary} 


\newtheorem{obs}[theorem]{Observation}
%\newtheorem{lemma}[theorem]{Lemma}
\newtheorem{claim}[theorem]{Claim}
\newtheorem*{statement}{Statement}
\newtheorem{corollary}[theorem]{Corollary}
\newtheorem{conj}[theorem]{Conjecture}
\newtheorem{lemma}[theorem]{Lemma}
%\newtheorem{thrm}[theorem]{Theorem}
% to make my stuff easier to write, i use \begin{defn} \end{defn} and here i define the shortcuts:
\theoremstyle{definition}
\newtheorem{alg}[theorem]{Algorithm}
\newtheorem{definition}[theorem]{Definition}
\newtheorem{question}[theorem]{Question}
\newtheorem{remark}[theorem]{Remark}
\newtheorem{notation}[theorem]{Notation}
\newtheorem*{assumption}{Assumption}
\newtheorem{example}[theorem]{Example}
\newtheorem*{ex*}{Example}
\newtheorem*{problem}{Problem}
\newtheorem{exercise}[theorem]{Exercise}

\newtheorem{defprop}[theorem]{Definition-Proposition}
\newtheorem*{Dthrm}{Decomposition Theorem}


%%solving the link does not correspond to correct numer
\makeatletter
\renewcommand*{\theHdefinition}{\thesection.\arabic{definition}}
\renewcommand*{\theHtheorem}{\thesection.\arabic{theorem}}
\renewcommand*{\theHlemma}{\thesection.\arabic{lemma}}
\renewcommand*{\theHproposition}{\thesection.\arabic{proposition}}
\renewcommand*{\theHequation}{\thesection.\arabic{equation}} 
\renewcommand*{\theHclaim}{\thesection.\arabic{claim}} 
\renewcommand*{\theHremark}{\thesection.\arabic{remark}}
\renewcommand*{\theHconj}{\thesection.\arabic{conj}}
\makeatother
%%


\usepackage[hyperpageref]{backref}
\renewcommand{\backref}[1]{$\uparrow$~#1}

\theoremstyle{plain}
\newlist{thmlist}{enumerate}{1}
\setlist[thmlist]{wide = 0pt, labelwidth = 2em, labelsep*=0em, itemindent = 0pt, leftmargin = \dimexpr\labelwidth + \labelsep\relax, noitemsep,topsep = 1ex, font=\normalfont, label=(\roman*), ref=\thetheorem.(\roman{thmlisti})}


\newlist{thmenum}{enumerate}{1} % also creates a counter called 'propenumi'
\setlist[thmenum]{wide = 0pt, labelwidth = 2em, labelsep*=0em, itemindent = 0pt, leftmargin = \dimexpr\labelwidth + \labelsep\relax, noitemsep,topsep = 1ex, font=\normalfont, label=(\roman*), ref=\thethmx.(\roman{thmenumi})}%{label=\alph*), ref=\thethmx~(\alph*)}
\crefalias{thmenumi}{thmx} 

\newlist{corlist}{enumerate}{1} % also creates a counter called 'propenumi'
\setlist[corlist]{wide = 0pt, labelwidth = 2em, labelsep*=0em, itemindent = 0pt, leftmargin = \dimexpr\labelwidth + \labelsep\relax, noitemsep,topsep = 1ex, font=\normalfont, label=(\roman*), ref=\thecorx.(\roman{corlisti})}%{label=\alph*), ref=\thethmx~(\alph*)}
\crefalias{corlisti}{corx} 

\addtotheorempostheadhook[theorem]{\crefalias{thmlisti}{theorem}}


\addtotheorempostheadhook[corollary]{\crefalias{thmlisti}{corollary}}

\addtotheorempostheadhook[proposition]{\crefalias{thmlisti}{proposition}}

\addtotheorempostheadhook[definition]{\crefalias{thmlisti}{definition}}

\addtotheorempostheadhook[lemma]{\crefalias{thmlisti}{lemma}} 



\addtotheorempostheadhook[remark]{\crefalias{thmlisti}{remark}}



\crefname{lemma}{Lemma}{Lemmas} 
\crefname{conjecture}{Conjecture}{Conjectures}
\crefname{theorem}{Theorem}{Theorems}
\crefname{proposition}{Proposition}{Propositions}
\crefname{definition}{Definition}{Definitions}
\crefname{remark}{Remark}{Remarks}
\crefname{corollary}{Corollary}{Corollaries} 
\crefname{corx}{Corollary}{Corollaries} 
\crefname{thmx}{Theorem}{Theorems}


%%solving the link does not correspond to correct numer
\makeatletter
\renewcommand*{\theHdefinition}{\thesection.\arabic{definition}}
\renewcommand*{\theHtheorem}{\thesection.\arabic{theorem}}
\renewcommand*{\theHlemma}{\thesection.\arabic{lemma}}
\renewcommand*{\theHproposition}{\thesection.\arabic{proposition}}
\renewcommand*{\theHequation}{\thesection.\arabic{equation}} 
\renewcommand*{\theHclaim}{\thesection.\arabic{claim}} 
\renewcommand*{\theHremark}{\thesection.\arabic{remark}}
%\renewcommand*{\theHconjecture}{\thesection.\arabic{conjecture}}
\makeatother
%%

\newcommand\bQ{{\mathbb Q}}

\newcommand\sA{{\mathcal A}}
\newcommand\sO{{\mathcal O}}
\newcommand\sK{{\mathcal K}}
\newcommand\sH{{\mathscr H}}
\newcommand\sC{{\mathscr C}}
\newcommand\sD{{\mathscr D}}

\newcommand\sP{{\mathcal P}}
\newcommand\cP{{\mathcal P}}
\newcommand\sI{{\mathcal I}}
\newcommand\sJ{{\mathcal J}}
\newcommand\sF{{\mathcal F}}
\newcommand\sE{{\mathcal E}}
\newcommand\sG{{\mathcal G}}
\newcommand\sM{{\mathcal M}}
\newcommand\sW{{\mathcal W}}
\newcommand\sV{{\mathcal V}}
\newcommand\sL{\mathcal{L}}
\def\cU{\mathcal{U}}
\newcommand\sR{\mathcal{R}}
\newcommand\sB{\mathscr{B}}
\newcommand\sS{\mathcal{S}}
\newcommand\OS{\mathrm{OS}}
\newcommand\tr{\mathrm{tr}}

\newcommand\GL{\mathrm{GL}}

\def\cA{\mathcal{A}}
\def\sQ{\mathcal{Q}}
\def\R{\mathbb{R}}
\def\K{\mathbb{K}}
\def\bA{\mathbb{A}}
\def\kC{\mathfrak{C}}
\def\N{\mathbb{N}}
\def\Z{\mathbb{Z}}
\def\bG{\mathbb{G}}
\def\bF{\mathbb{F}}
\def\bE{\mathbb{E}}

\def\bV{\mathbb{V}}
\def\bM{\mathbb{M}}
\def\ep{\varepsilon}
\def\brho{{\bm{\varrho}}}

\def\btau{{\bm{\tau}}}

\def\bsigma{{\bm{\sigma}}}

\newcommand{\Set}{\mathsf{SET}}
\newcommand{\sAaug}{\mathcal{A}^{\bullet}_{\textrm{DR}, 0}}
\newcommand{\ltensor}{\stackrel{L}{\otimes}}
\newcommand{\RHOM}{\mathbf{R}\mathcal{H}om}

\newcommand{\vhs}{{\rm VHS}}

\newcommand\an{\textup{an}}


%\DeclareMathOperator{\Pic}{Pic}                  % Pic

%\renewcommand{\char}{\mathrm{char}}           % Poincare line bunle

\DeclareMathOperator{\Supp}{Supp}                % Supp
\DeclareMathOperator{\codim}{codim}              % codim
\DeclareMathOperator{\mreg}{mreg}                % mreg
\DeclareMathOperator{\reg}{reg}                  % reg
\DeclareMathOperator{\sing}{sing}                  
\DeclareMathOperator{\id}{id}                    % id
\DeclareMathOperator{\obj}{Obj}
\DeclareMathOperator{\ad}{ad}
\DeclareMathOperator{\morph}{Morph}
\DeclareMathOperator{\enom}{End}
\DeclareMathOperator{\iso}{Iso}
\DeclareMathOperator{\Exp}{Exp}
\DeclareMathOperator{\homo}{Hom}
\DeclareMathOperator{\enmo}{End}
\DeclareMathOperator{\norm}{norm}
\DeclareMathOperator{\spec}{Spec}
\DeclareMathOperator{\univ}{univ}
\DeclareMathOperator{\mspec}{mSpec}
\DeclareMathOperator{\fitt}{Fitt}
\DeclareMathOperator{\odr}{\Omega^\bullet_{\textrm{DR}}}
\DeclareMathOperator{\rank}{Rank}
\def\Nat{{\rm Nat}}
\DeclareMathOperator{\grad}{grad}
\DeclareMathOperator{\Crit}{Crit}
\def\hess{\sqrt{-1}\partial\bar{\partial}}
\def\ddc{{dd^{\mathrm{c}}}}
%\DeclareMathOperator{\char}{char}
\def\db{\bar{\partial}}


%\DeclareMathOperator

\DeclareMathOperator{\red}{red}
\DeclareMathOperator{\quot}{Quot}
\DeclareMathOperator{\kn}{Ker}
\DeclareMathOperator{\im}{Im}
\DeclareMathOperator{\dvol}{dvol}

\def\bone{\mathbf{1}}
\def\C{\mathbb{C}}
\def\bR{\mathbb{R}}
\def\cM{{\mathscr{M}}}
\def\cV{\mathcal{V}}
\def\Def{{\rm {Def}}}
\def\cR{\mathscr{R}}
\def\cB{\mathcal{B}}

\def\om{\omega}
\def\wti{\widetilde}
\def\al{\alpha}
\def\End{{\rm {End}}}
\def\Pic{{\rm Pic}}
\def\cE{\mathcal{E}}
\def\cF{\mathcal{F}}
\def\bC{\mathbb{C}}
\def\bP{\mathbb{P}}
\def\cH{\mathcal{H}}
\def\bL{\mathbb{L}}
\def\cX{\mathcal{X}}
\def\cI{\mathcal{I}}
\def\pa{\partial}
\def\cY{\mathcal{Y}}
\def\cC{\mathcal{C}}
\def\cD{\mathscr{D}}
\def\cO{\mathcal{O}}
\def\lra{\longrightarrow}
\def\Q{\mathbb{Q}}
\def\cL{\mathcal{L}}
\def\cG{\mathcal{G}}
\def\bD{\mathbf{D}}
\def\cS{\mathcal{S}}

\def\kg{\mathfrak{g}}
\def\ku{\mathfrak{u}}

\def\xsp{X^{\rm sp}}
\def\ysp{Y^{\rm sp}}

\def\zsp{Z^{\rm sp}}




\author[Y. Deng]{Ya Deng}
\address{CNRS, Institut \'Elie Cartan de Lorraine, Universit\'e de Lorraine, Site de
	Nancy,   54506 Vand\oe uvre-lès-Nancy, France} 
\email{ya.deng@math.cnrs.fr}
\urladdr{https://ydeng.perso.math.cnrs.fr} 

\author[B. Wang]{Botong Wang}
\address{Department of Mathematics, University of Wisconsin-Madison, 480 Lincoln Drive, Madison WI 53706-1388, USA}
\email {wang@math.wisc.edu}
\urladdr{https://people.math.wisc.edu/~bwang274/}
 
%\title[Vanishing theorems for $L^2$-cohomology]{Linear Koll\'ar Conjecture on holomorphic Euler characteristic}
\title[ $L^2$-vanishing theorem and a conjecture of Koll\'ar]{$L^2$-vanishing theorem 
	and a conjecture of Koll\'ar}
\alttitle{Théorème d’annulation \(L^2\) et une conjecture de Kollár}
\keywords{Koll\'ar conjecture,  $L^2$-cohomology, Vanishing theorem, Harmonic map, Bruhat-Tits  building, Linear Shafarevich conjecture, Variation of (mixed) Hodge structure, Generically large fundamental group}
\altkeywords{Conjecture de Kollár, Cohomologie \(L^2\), Théorème d’annulation, Application harmonique, Immeuble de Bruhat–Tits, Conjecture de Shafarevich linéaire, Variation de structure de Hodge (mixte), Groupe fondamental génériquement grand}
\subjclass{32J25, 14D07, 53C23, 32L20}

\begin{document} 
	
	\begin{abstract}
		%		Let $X$ be a smooth complex projective variety of dimension $n$ and let $\varrho:\pi_1(X)\to {\rm GL}_N(\bC)$ be a linear representation. In this paper   we mainly prove that
	In 1995, Koll\'ar conjectured that   a smooth complex projective \(n\)-fold \(X\) with generically large fundamental group has  Euler characteristic \(\chi(X, K_X)\geq 0\).      In this paper, we confirm  the conjecture assuming \(X\) has  linear fundamental group, i.e., there exists a  representation \(\pi_1(X)\to {\rm GL}_N(\mathbb{C})\) with finite kernel.  We deduce the conjecture by proving a stronger \(L^2\) vanishing theorem: for the universal cover \(\widetilde{X}\) of such \(X\), its   \(L^2\)-Dolbeault cohomology \(H_{(2)}^{n,q}(\widetilde{X})=0\) for  \(q\neq 0\).     The main ingredients of the proof are techniques from the linear Shafarevich conjecture along with  some analytic methods. 


%As a byproduct, we  give geometric characterization of having generically large fundamental group. 

		%we prove that if $\pi_1(X)$ is linear,  i.e., there exists an almost faithful   representation  $\pi_1(X)\to {\rm GL}_N(\bC)$,  then $\pi_1(X)$ is large (resp. generically large) if and only if  $X$ is K\"ahler parabolic (resp. semi-K\"ahler parabolic). This gives an analytic characterization of varieties with (generically) large fundamental groups. To  prove the above theorems, we substantially apply the techniques from the proofs of the reductive and linear Shafarevich conjectures, along with  $L^2$-vanishing theorems and   the geometry of mixed period maps for variations of mixed Hodge structures.   %We also proved that $\chi(X,K_X)\geq 0$ if there exists a large and semisimple representation  $\varrho:\pi_1(X)\to \GL_N(\bC)$, which addresses a conjecture by Koll\'ar. 
	\end{abstract}
	\begin{altabstract}
		En 1995, Kollár a conjecturé qu’une variété projective complexe lisse \(X\) de dimension \(n\), dont le groupe fondamental est génériquement  grand, vérifie l’inégalité \( \chi(X, K_X) \geq 0 \). Dans cet article, nous confirmons cette conjecture sous l’hypothèse que le groupe fondamental de \(X\) est linéaire, c’est-à-dire qu’il existe une représentation \( \pi_1(X) \to \mathrm{GL}_N(\mathbb{C}) \) à noyau fini. Notre approche repose sur un théorème d’annulation \(L^2\) plus fort: pour le revêtement universel \( \widetilde{X} \) de \( X \), la cohomologie de Dolbeault \(L^2\) vérifie \( H_{(2)}^{n,q}(\widetilde{X}) = 0 \) pour tout \( q \neq 0 \). Les ingrédients principaux de la démonstration proviennent de techniques liées à la conjecture de Shafarevich linéaire, ainsi que de méthodes analytiques.
	\end{altabstract}
	\maketitle
	\tableofcontents
	\section{Introduction}
	\subsection{Main theorem}
	In the  study of Shafarevich maps, Koll\'ar made the following conjecture (\cite[Conjecture 18.12.1]{Kol95}).
\begin{conj}[Koll\'ar]\label{conj:kollar}
	Let $X$ be a smooth complex projective variety. If $X$ has generically large fundamental group, then $\chi(X,K_X)\geq 0$. 
\end{conj} 
	Following the notations of \cite{Kol95}, we say that $X$ \emph{has generically large fundamental group} (resp. $\varrho:\pi_1(X)\to \GL_N(\bC)$ is a \emph{generically large representation}) if for any irreducible positive-dimensional subvariety $Z$ of $X$ passing through a very general point, the image ${\rm Im}[\pi_1(Z^{\rm norm})\to \pi_1(X)]$ (resp. $\rho({\rm Im}[\pi_1(Z^{\rm norm})\to \pi_1(X)])$) is infinite (see \cite[Definition 2.4]{Kol95} for the precise meaning of ``very general''). Here $Z^{\rm norm}\to Z$ is the normalization of $Z$.     Note that in \cite{CDY22,DY23,DY24}, generically large is called \emph{big}.  
	
In \cite{GL87}, Green and Lazarsfeld  proved Koll\'ar's conjecture when $X$ has maximal Albanese dimension.   In this paper, we will use  methods of $L^2$-cohomology to study  \Cref{conj:kollar}. For a compact  K\"ahler manifold $(X,\omega)$, we denote by $\pi_X:\widetilde{X}\to X$ its universal cover, and  for any non-negative integer $p$ and $q$,  let  $H_{(2)}^{p,q}(\widetilde{X})$ be  the  $L^2$-Dolbeault cohomology group  with respect to the metric $\pi_X^*\omega$. Note that its  definition does not depend on the choice of $\omega$.
	Our main theorem is the following:	\begin{thmx}\label{main:kollar}
		Let $X$ be a  smooth complex projective  $n$-fold. If there exists a generically large representation $\varrho:\pi_1(X)\to \GL_N(\bC)$, then the following statements hold.
		    \begin{thmenum} 
			\item  \label{item:vanishing} $H^{p,0}_{(2)}(\widetilde{X})=0$   for  $0\leq p\leq n-1$ and $H^{n,q}_{(2)}(\widetilde{X})=0$ for  $1\leq q\leq n$. 
			\item \label{item:Euler}The Euler characteristic $\chi(X, K_X)\geq 0$.
			\item  \label{item:strict} If the strict inequality $\chi(X,K_X)>0$ holds, then 
			\begin{enumerate}[label*=(\alph*)]
				\item \label{item:converse2}there exists  {a nontrivial} $L^2$-holomorphic $n$-form on $\widetilde{X}$;
				\item  \label{item:converse} $X$ is of general type. 
			\end{enumerate} 
		\end{thmenum}
	\end{thmx}
	In particular, we prove \Cref{conj:kollar} assuming $\pi_1(X)$ is linear, i.e. there exists a representation $ \pi_1(X)\to \GL_N(\bC)$ with finite kernel. 
	
	The most difficult aspect of proving \cref{main:kollar} is  showing that $H^{p,0}_{(2)}(\widetilde{X})=0$   for  $0\leq p\leq n-1$. We  outline the proof strategy at the beginning of  \Cref{sec:linear}. The  conclusion of \cref{main:kollar} follows from this $L^2$-vanishing result via Atiyah's $L^2$-index theorem. This  implication from \cref{item:vanishing} to \cref{item:Euler} is well known to experts, and similar results have  appeared in various forms in the literature, including \cite{Gr,CD01,Eys99,Eys00,Tak04,JZ}, among others.
	
Let us point out that in the special case where $\varrho$ in \cref{main:kollar} is assumed to be \emph{semisimple}, both \cref{item:vanishing} and \cref{item:Euler} were previously claimed in \cite{JZ}, although certain aspects of their argument appear to require further clarification. Our approach follows a similar overall strategy; however, the novelty in the linear case lies in establishing a partial vanishing theorem  (see \cref{thm:vanishing2}), which requires more refined techniques from harmonic maps into Euclidean buildings  in order to apply the methods developed in the linear Shafarevich conjecture \cite{EKPR12} and thereby prove \cref{item:vanishing}.
	\iffalse
	\subsection{Strategy of the proof}
	We first introduce the notion of semi-K\"ahler parabolicity in \cref{def:wkp}. We then prove the following $L^2$-vanishing theorem, generalizing previous works \cite{Gr,Esy97,CX,JZ,BDET22}. % for  (semi-)K\"ahler parabolic manifolds.  
	\begin{thmx}[=\cref{thm:vanishing2}]\label{main3}
		Let $(X,\omega)$ be a   compact  semi-K\"ahler parabolic manifold of dimension $n$.  Then 	we have \begin{thmenum}[resume]
			\item \label{item:vanishing2}$H^{p,0}_{(2)}(\widetilde{X})=0$, \;for\;  $0\leq p\leq n-1$;
			\item \label{item:L2} $H^{n,q}_{(2)}(\widetilde{X})=0$, \;for\; $1\leq q\leq n$. 
		\end{thmenum}
	\end{thmx}
	Next, using techniques from the reductive and linear Shafarevich conjectures in \cite{DY23,EKPR12} as well as tools in the harmonic maps to Euclidean buildings in \cite{GS92,BDDM} and geometry of mixed period maps,% of variation of mixed Hodge structures, 
	we prove the following result.
	\begin{thmx}[=\cref{thm:main2}]\label{main}
		Let $X$ be a smooth projective variety of dimension $n$. Let $\varrho:\pi_1(X)\to \GL_N(\bC)$ be a linear representation. 
	 If $\varrho$ is  generically large,  then  $X$ is semi-K\"ahler parabolic. %In particular, if  $\pi_1(X)$ linear, then  $X$ is      semi-K\"ahler parabolic   if and only if $X$ has generically large fundamental group. 
	\end{thmx}   
	Let us explain how to derive \cref{main:kollar} from \cref{main,main3}.
	 
	\begin{proof}[Proof of \cref{main:kollar}]
		By \cref{main:kollar}, $X$ is semi-K\"ahler parabolic. 	The first claim follows from \cref{main3}.	
		
		We denote by $\Gamma=\pi_1(X)$ and $\dim_{\Gamma}H^{n,q}_{(2)}(\widetilde{X})$ the Von Neumann dimension of $H^{n,q}_{(2)}(\widetilde{X})$ (cf. \cite{Ati76} for the definition). 
		By Atiyah's $L^2$-index theorem  and \Cref{item:L2}, we have
		\begin{align} 	\chi(X,K_X)=\sum_{q=0}^{n}(-1)^q\dim_{\Gamma}H^{n,q}_{(2)}(\widetilde{X})= \dim_{\Gamma}H^{n,0}_{(2)}(\widetilde{X})\geq 0.
		\end{align} 
 The second claim is proved. 
	
		
		If the strict  inequality in \eqref{eq:euler} holds, then $H^{n,0}_{(2)}(\widetilde{X})\neq 0$. We then apply \cite[Corollary 13.10]{Kol95} to conclude that $K_X$ is big. The third claim follows. 
	\end{proof} 
	For the proof of \cref{main:kollar}, when $\varrho$ is semisimple, 	we shall construct  suitable  1-forms required in \cref{def:wkp} via the harmonic map to Euclidean buildings induced by some representations $\pi_1(X)\to \GL_N(K)$ with $K$ a non-archimedean local field. %, as derived from the proof of the reductive Shafarevich conjecture \cite{DY23}. 
	If $\varrho$ is merely linear, the construction of these 1-forms are quite involved, derived from the geometry of mixed period domains together with the structure of  the linear Shafarevich morphism.  While we could bypass the semisimple case and directly prove \cref{main:kollar} for linear representation, we choose to write the proof of the semisimple case to isolate the more technical arguments and benefits readers who are more interested in the semisimple case.  
	\fi
	%	Another conjecture by Koll\'ar \cite[Conjecture 18.8.1]{Kol95}  predicts that if $\pi_1(X)$ is generically large  and $X$ is of general type, then $\chi(X,K_X)>0$ and there exist nontrivial $L^2$-holomorphic $n$-forms on $\widetilde{X}$. One might expect that for  $X$ as in \cref{main:kollar},  this conjecture holds. However, in \cite{EL97}  Ein and Lazarsfeld   constructed a smooth projective threefold $X$ with maximal Albanese dimension (hence $\pi_1(X)$ has a generically large and abelian representation) such that $X$ is of general type with $\chi(X,K_X)=0$.  Therefore, by the proof of \cref{main:kollar}, $\widetilde{X}$ does not admit any $L^2$-holomorphic $3$-form. More examples can be found in \cite{CDJ14}.
	
	
	\subsection{Some histories and comparison with previous works}
	In this subsection, $(X,\omega)$ is a compact K\"ahler manifold of dimension $n$ and we denote by $\pi_X:\widetilde{X}\to X$  the universal cover.  In \cite{Gr}, Gromov introduced the notion of \emph{K\"ahler hyperbolicity} in his study of the Hopf conjecture. Recall that  $X$ is K\"ahler hyperbolic if there is a smooth 1-form $\beta$ such that $\pi_X^*\omega=d\beta$ and the norm $|\beta|_{\pi_X^*\omega}$ is bounded from above by a constant.   He then proved the vanishing of $k$-th $L^2$-Betti numbers of $\widetilde{X}$  for $k\neq n$. 
	
	
	Gromov's idea was later extended by Eyssidieux \cite{Esy97}, in which  more general notions  of \emph{weakly K\"ahler hyperbolicity} are introduced.   Eyssidieux's  work was recently generalized by {Bei}, Claudon, {Diverio}, {Eyssidieux}, and {Trapani} \cite{BDET22,BCDT24}, who studied a {birational} analog of {K\"ahler hyperbolicity}. 
 Also as generalizations of Gromov's work,   Cao-Xavier \cite{CX} and Jost-Zuo \cite{JZ}   introduced the notion of \emph{K\"ahler parabolicity}. They independently observed that the arguments of Gromov concerned with the vanishing of $L^2$-Betti numbers work also under the weaker assumption that  $\beta$ is a smooth 1-form such that $|\beta(x)|_{\pi_X^*\omega}$  has \emph{sub-linear growth}. In other words,   there exists a constant $c>0$ such that  $|\beta(x)|_{\pi_X^*\omega}\leq c(1+d_{\widetilde{X}}(x,x_0))$, where $d_{\widetilde{X}}(x,x_0)$ is the Riemannian distance between $x$ and the base point $x_0$ in $\widetilde{X}$ relative to the metric $\pi^*\omega$.   %Recall that  $X$ is \emph{K\"ahler parabolic} if $\pi_X^*\omega=d\beta$, where  $\beta$ a smooth 1-form such that $|\beta(x)|_{\pi_X^*\omega}$  has \emph{sub-linear growth}. In other words,   there exists a constant $c>0$ such that  $|\beta(x)|_{\pi_X^*\omega}\leq c(1+d_{\widetilde{X}}(x,x_0))$, where $d_{\widetilde{X}}(\bullet,\bullet)$ is the Riemannian distance between $x$ and the base point $x_0$ in $\widetilde{X}$ relative to the metric $\pi^*\omega$.  
	
The formulation and proof of our partial $L^2$-vanishing theorem in  \cref{thm:vanishing2}  are inspired by the aforementioned works, and is introduced specifically for the proof of \cref{main:kollar}.  As a result, it may appear technically involved.  
	
Another key component of the proof of \cref{main:kollar} draws on techniques developed in the context of the reductive and linear Shafarevich conjectures, as explored  in \cite{DY23, Eys04, EKPR12}. The Shafarevich conjecture predicts that the universal cover of a smooth complex projective variety is holomorphically convex. This was proved by Eyssidieux, Katzarkov, Pantev, and Ramachandran \cite{Eys04, EKPR12} for smooth projective varieties with linear fundamental groups. More recently, the conjecture was extended by Yamanoi and the first author in \cite{DY23} to projective normal varieties with reductive fundamental groups, and further studied in \cite{DY24} for projective normal varieties with fundamental groups in positive characteristic.
	
	
	
	
	\subsection{Notation and Convention}
	\begin{itemize}
	\item All varieties in this paper are defined over $\C$.
		\item Let $(X.\omega)$ be a K\"ahler manifold. We denote by $\pi_X:\widetilde{X}\to X$ the universal cover of $X$.  
		\item Let $(X,\omega)$ be a compact K\"ahler manifold. Unless otherwise specified, $d_{\widetilde{X}}(x, x_0)$ stands for the Riemannian distance between $x$ and the base point $x_0$ in $\widetilde{X}$ with respect to the metric $\pi_X^*\omega$. %, and  $H_{(2)}^k(\widetilde{X},\bC)$ denotes the $L^2$-cohomology with respect to the $\pi_X^*\omega$, which does not depend the choice of $\omega$.  
		\item For a complex space $Z$, $Z^{\rm norm}$ denotes its normalization, and $Z^{\rm reg}$ denotes its regular locus. 
		\item We use the standard abbreviations VHS and VMHS for \emph{variation of Hodge structure} and \emph{variation of mixed Hodge structure} respectively.
		\item By convention, a closed positive $(1,1)$-current $T$ on a complex manifold is   \emph{semi-positive}. 
		\item Plurisubharmonic functions  are  abbreviated as psh functions. 
		\item A positive closed $(1,1)$-current $T$  has \emph{continuous potential} if locally we have $T=\hess\psi$ with $\psi$  a continuous psh function. 
		\item By convention, a real \emph{positive} \( (1,1) \)-form or current is understood to be \emph{semi-positive} in the usual sense. 
	\end{itemize}
	\subsection{Acknowledgment}    We would like to express our gratitude to Patrick Brosnan, Junyan Cao, Simone Diverio, Philippe Eyssidieux, Mihai P\u{a}un, Pierre Py, Christian Schnell,  Xu Wang, and Mingchen Xia for their inspiring discussions. We extend special thanks to Chikako Mese for her invaluable discussion  and expertise on harmonic maps into Euclidean buildings, and  to Yuan Liu and Xueyuan Wan for reading the first version of the paper and their very helpful remarks.   YD acknowledges the partial support of the French Agence Nationale de la Recherche (ANR) under reference ANR-21-CE40-0010. BW thanks Peking University and BICMR for the generous hospitality and support during the writing of this paper. He is partially supported by a Simons fellowship.
	
	\section{A partial $L^2$-vanishing theorem}
 \label{sec:kp}
%A closed positive $(1,1)$-current $T$ on a complex manifold $X$ can be seen as a $(1,1)$-form with  positive measure coefficients. Note that positive measures  admit
%a Lebesgue decomposition into an absolutely continuous part (with respect to the
%Lebesgue measure on $X$) and a singular part. We therefore get a decomposition
%of $T$ itself into an absolutely continuous part $T_{\rm ac}$ and a singular part $T_{\rm sing}$.  We begin with the following definition in \cite[\S 2.3]{Bou02}. \begin{definition}[Lebesgue decomposition]\label{def_Leb}
%The
%	decomposition $T=T_{\rm ac}+T_{\rm sing}$ is called the \emph{Lebesgue decomposition} of $T$ . 
%	\end{definition}
%		\begin{definition}[Semi-K\"ahler form]
%		Let $(X,\omega_X)$ be a compact K\"ahler manifold. A  smooth closed real $(1,1)$-form $\omega_{\rm sk}$ on  $X$ is \emph{semi-K\"ahler} if $\omega_{\rm sk}$ is semipositive everywhere and  is strictly positive on a Zariski dense open subset of $X$. 
%	\end{definition} 
	
 
	%\begin{remark}  Note that \( r(x) := d_{\widetilde{X}}(x, x_0) + 1 \) is continuous, exhausting, and smooth almost everywhere. Moreover, there exists a constant \( c > 0 \) such that \( |dr| \leq_{\text{a.e.}} c \). Hence, it fulfills the conditions in \Cref{item:dlinear}. We mention here that our definition of K\"ahler parabolicity is much weaker than that in \cite{CX}. The definition of K\"ahler parabolic manifolds by Cao-Xavier in \cite{CX} corresponds to our strongly K\"ahler parabolic ones with \( r(x) = d_{\widetilde{X}}(x, x_0) + 1 \). The relaxation of conditions for \( r \) is crucial in the proof of \cref{thm:main2}. Indeed, in \cref{thm:main}, we will use the function \( d_{\widetilde{X}}(x, x_0) + 1 \) to prove (semi)-K\"ahler parabolicity for semisimple representations of $\pi_1(X)$. However, if the representation is merely linear, we can   construct a smooth 1-form \( \beta \) satisfying \Cref{item:bounded}, which merely satisfies \Cref{item:dlinear} with respect to another suitable exhausting function \( r \) instead of \( d_{\widetilde{X}}(x, x_0) + 1 \). 
	%\end{remark}
	
	
	
	
	Before proving a key $L^2$-vanishing theorem, we first recall the definition of $L^2$-cohomology (cf. \cite[Definition 12.3]{DIP02}).
	\begin{definition}[$L^2$-cohomology]\label{defn:L2} 
		Let $(Y,\omega)$ be a complete K\"ahler manifold.  Let   $L_{(2)}^{p,q}(Y)$  be the space of $L^2$-integrable   $(p,q)$-forms with respect to the metric $\omega$. A section $u$ is said to be in  ${\rm Dom}\, \db$  if  $\db u$  calculated in the sense of distributions is still in $L^2$.  Then  the $L^2$-Dolbeault cohomology is defined as
		\[ H^{p,q}_{(2)}(Y)= {{\rm ker}\, \db}\big/\,{\overline{{\rm Im}\, \db \cap {\rm Dom}\, \db}}.
		\]  
		If $(X,\omega)$ is a compact K\"ahler manifold, then 	$H^{p,q}_{(2)}(\widetilde{X})$ denotes the $L^2$-cohomology with respect to the metric $\pi_X^*\omega$. It is well known that this cohomology group is independent of the choice of Kähler metric $\omega$. 
		%		$
		%		H^k_{(2)}(X,\bC)= {{\rm ker}\, d}/{\overline{{\rm Im}\, d\cap {\rm Dom}\, d}}
		%		$ (resp.  $
		%		H^{p,q}_{(2)}(X)= {{\rm ker}\, \db}/{\overline{{\rm Im}\, \db \cap {\rm Dom}\, \db}}
		%		$).
	\end{definition} 
	
	Let us state and prove our partial  $L^2$-vanishing theorem. 
	\begin{theorem} \label{thm:vanishing2} 
 Let   $(X,\omega)$ be a compact 	 K\"ahler $n$-fold.  Let $f:X\to Y$ be a proper surjective holomorphic map with connected fibers over a compact  K\"ahler normal space $Y$ of dimension $m$. Denote by $\tilde{f}:\widetilde{X}\to \widetilde{Y}$ the lift of $f$ to  universal covers.   Assume that  there exists   a  1-form  $\beta$ on $\widetilde{X}$  with $L_{\rm loc}^1$-coefficients  satisfying the following properties. 
 	\begin{enumerate}[label*=(\rm \alph*)]
 	\item \label{item:smooth} The 1-form $\beta$ is smooth  on an open subset   of $\widetilde{X}$ whose complement has zero Lebesgue measure, and  there is a constant $C>0$ such that  \begin{align}\label{item:dlinear}
 	 |\beta(x)|_{\pi_X^*\omega}\leq_{\rm a.e.} C(d_{\widetilde{X}}(x,x_0)+1)
 \end{align}
 where $x_0$ is a base point of $\widetilde{X}$. 
 	\item\label{item:bounded}   The current \( d\beta \) is a real positive \( (1,1) \)-current satisfying
 	\begin{align}\label{eq:mutual}
 		d\beta \geq \pi_X^* f^* T_Y,
 	\end{align}
 	where \( T_Y \) is a closed positive \( (1,1) \)-current on \( Y \). Moreover, there exists a non-empty analytic open subset \( Y_0\subset Y \) such that \( T_Y|_{Y_0} \) is smooth and strictly positive.
 \end{enumerate} 
 Then  \begin{thmlist}
 	 \item\label{item:partial}  For any $p\in \{0,\ldots,m-1\}$, we have  $H^{p,0}_{(2)}(\widetilde{X})=0$. 
 	\end{thmlist}
Assume that there exists a non-zero class $\alpha \in H^{p,0}_{(2)}(\widetilde{X})$ for some $p\in \{m,\ldots,n-1\}$. Let $Y^\circ$ be a non-empty open subset of $Y^{\mathrm{reg}} \cap Y_0$ over which $f$ is a proper submersion, and such that the restriction $T_Y|_{Y^\circ}$ is a K\"ahler form.  Then
 	\begin{thmlist} [resume]
 	 \item\label{item:long}  we have   \begin{align}\label{eq:tensor}
 	 	 \alpha|_{\widetilde{X}^\circ}\in H^0(\widetilde{X}^\circ, \tilde{f}^*\Omega_{\widetilde{Y}^\circ}^{m}\otimes \Omega_{\widetilde{X}^\circ}^{p-m}), 
 	 \end{align}
where $\widetilde{Y}^\circ:= \pi_Y^{-1}(Y^\circ)$ and $\widetilde{X}^\circ:=\tilde{f}^{-1}(\widetilde{Y}^\circ)$.
\item \label{item:d-closed} For any $y\in \widetilde{Y}^\circ$,  let $\alpha_y$ be the holomorphic $(p-m)$-form on $\tilde{f}^{-1}(y)$ induced by $\alpha$   under the isomorphism $$\tilde{f}^*\Omega_{\widetilde{Y}^\circ}^{m}\otimes \Omega_{\widetilde{X}^\circ}^{p-m}|_{\tilde{f}^{-1}(y)}\simeq \Omega_{\tilde{f}^{-1}(y)}^{p-m}.$$ 
 	Then  $
 	 \alpha|_{\tilde{f}^{-1}(y)}
 	 $ is   $d$-closed.
 \end{thmlist}
	
	\end{theorem}
	\begin{proof}
  To lighten the notation, we write  $\omega$   instead of   $\pi_X^*\omega$ abusively.  
   
   \noindent {\bf Step 1. }
Since the sectional curvature of the complete K\"ahler manifold $(\widetilde{X},\pi_X^*\omega)$ is uniformly bounded,  by a result of W. Shi (see e.g. \cite[Theorem 1.2]{Hua19}),  there exists a constant $C>0$ and a  smooth exhausting function $r$   on $\widetilde{X}$ such that 
$$
d_{\widetilde{X}}(x,x_0)+1 \leq r(x)\leq d_{\widetilde{X}}(x,x_0)+C,
$$ 
and $|d r|_{\omega}(x) \leq C$  for all $x \in \widetilde{X}$.   Hence by \eqref{item:dlinear}, we have 
\begin{align}\label{eq:betanorm}
	 |\beta|_\omega(x)\leq_{\rm a.e.} Cr(x).
\end{align} 	Let $\varrho: \mathbb{R} \rightarrow \mathbb{R}$ be a smooth function with $0 \leq \varrho \leq 1$ such that
		$$
\varrho(t)= \begin{cases}1, & \text { if } t \leq 0; \\ 0, & \text { if } t \geq 1.\end{cases}
		$$
		We consider the compactly supported function
		\begin{align}\label{eq:support}
			f_j(x)=\varrho(r(x)-{j}+1),
		\end{align} 
		where $j$ is a positive integer.  Then  $${\rm Supp}(f_j)\subset \{x\in \widetilde{X}\mid   r(x)\leq  {j} \},$$ 
		$$
		df_j(x)=\varrho'(r(x)-j+1) dr.
		$$ 
	Then  there exists some constant $c_1>0$ such that 
		\begin{align} \label{eq:uniform}
 |df_j(x)|_{\omega}\leq c_1. 
		\end{align}
		for any $x\in \widetilde{X}$. 
		
	%	  Denote by $n=\dim X$.  Let $\alpha$ be a  holomorphic $(p,0)$-form which is $L^2$ with respect to $\omega_X$ with $1\leq p\leq n-1$.  We define  $$\Theta=i^{p^2}\alpha\wedge\overline{\alpha}\wedge T^{n-p-1}.$$
	%	It is a  closed positive $(n-1,n-1)$-current, where $T=\ddc \phi$ and $\phi$ is a continuous psh function on $\widetilde{X}$.    Note that
	%	\begin{align*}
	%	0&=	\int_{\widetilde{X} }  d(f_j  d^c \phi\wedge\Theta) \\
	%	&	=   	\int_{\widetilde{X} }  df_j \wedge  d^c \phi\wedge\Theta +	\int_{\widetilde{X} }   f_j \wedge  \ddc \phi\wedge\Theta\\
	%	&=	\int_{\widetilde{X} }  df_j \wedge  d^c \phi\wedge\Theta +	\int_{\widetilde{X} }   i^{p^2}f_j\alpha\wedge \bar{\alpha}\wedge T^{n-p}
%	\end{align*} 
	
%	Note that $d^c\phi\wedge\Theta:=d^c(\phi \Theta)$. 
		
	  Let $\alpha$ be a  holomorphic $(p,0)$-form which is $L^2$ with respect to $\omega$ with $0\leq p\leq n-1$. Then 
	\begin{align}\label{eq:finite}
 \sqrt{-1}^{p^2}\int_{\widetilde{X}}\alpha\wedge \bar{\alpha}\wedge \omega^{n-p}	=	\frac{(n-p)!}{n!}\int_{\widetilde{X}  }    |\alpha|_{\omega}\omega^n <\infty.
		\end{align}
	
	
	 
		Denote by 
		$$
		B_j:=\{x\in \widetilde{X}\mid   j-1\leq r(x)\leq  j \}.
		$$  
	In what follows, 	for any smooth form $\gamma$ on $\widetilde{X}$, we denote by $|\gamma|$ its norm with respect to the metric $\omega$.  	Denote $\dvol:=\frac{\omega^n}{n!}$.    
Recall that  $|\beta(x)| \leq_{\rm a.e.} Cr(x)$ for some constant $C>0$. Hence by $\operatorname{supp}( df_j) \subset {B_{j}} $,  one has
		\begin{align} \nonumber
			\left|  \int_{\widetilde{X} }  df_j\wedge \beta\wedge\alpha\wedge \bar{\alpha}\wedge \omega^{n-p-1}  \right| 
			&\leq   \int_{  B_j} \left|  df_j\wedge \beta\wedge\alpha\wedge \bar{\alpha}\wedge \omega^{n-p-1}\right| \dvol   \\ \nonumber
			&\leq \int_{B_j}     |d f_j| \;  |\beta| \; | {\alpha}|^2 \; |\omega^{n-p-1}|   \dvol\\\label{eq:small}
			&\stackrel{\eqref{eq:uniform}\& \eqref{eq:betanorm} }{\leq} c_1Cj\int_{B_j}     | {\alpha}|^2   \dvol.  
		\end{align}
		
		\begin{claim}\label{claim:converge}
			There exists a subsequence $\left\{j_i\right\}_{i \geq 1}$ such that
			\begin{align}\label{eq:crucial}
				\lim _{i \rightarrow \infty} j_i  \int_{B_{j_i}}|\alpha(x)|^2 \dvol=0 .
			\end{align}  
		\end{claim}	 
		\begin{proof}
			If not, then there exists a positive constant $c'$ and a positive integer $n_0$ such that 
			$$
		j\int_{B_{j} }|\alpha(x)|^2 \dvol \geq c'>0
			$$
			for any $j\geq n_0$.
			This yields
			$$
			\begin{aligned}
			 \int_{\widetilde{X}}|\alpha(x)|^2 \dvol & \geq \sum_{j=n_0}^{\infty} \int_{B_{j}}  |\alpha(x)|^2 \dvol \\
				& \geq c' \sum_{j=n_0}^{+\infty} \frac{1}{j}=+\infty,
			\end{aligned}
			$$
		thereby yielding a contradiction, as $\alpha$ is assumed to be an $L^2$-form with respect to $\omega$.
		\end{proof}	
			\Cref{claim:converge} implies that there exists a subsequence $\left\{j_i\right\}_{i \geq 1}$ for which \eqref{eq:crucial} holds.  By the Stokes formula, we have
		\begin{align}\label{eq:essential} 
			\sqrt{-1}^{p^2}	\int_{\widetilde{X} } f_j  d\beta \wedge \alpha\wedge \bar{\alpha}\wedge \omega^{n-p-1}
			= \sqrt{-1}^{p^2}	\int_{\widetilde{X} }  (-df_j\wedge \beta)\wedge\alpha\wedge \bar{\alpha}\wedge \omega^{n-p-1}    
		\end{align} 
	 	By \cite[Chapter III, Example 1.2]{Dembook}, $\sqrt{-1}^{p^2}	  \alpha\wedge \bar{\alpha}$ is a positive $(p,p)$-form in the sense of \cite[Chapter III, Definition 1.1]{Dembook}. 	 \cite[Chapter III, Corollary 1.9 \& Proposition 1.11]{Dembook} imply that $\sqrt{-1}^{p^2}\alpha\wedge \bar{\alpha}\wedge \omega^{n-p-1}$ is a positive $(n-1,n-1)$-form, hence \emph{strongly positive} (not strictly positive!) in the sense of  \cite[Chapter III, Definition 1.1]{Dembook}. 	  Since $d\beta$ is assumed to be a positive current, it follows that
		 $$
 	\sqrt{-1}^{p^2}	   d\beta \wedge \alpha\wedge \bar{\alpha}\wedge \omega^{n-p-1}
		 $$
		 is a  positive measure on $\widetilde{X}$.   \eqref{eq:small}    and \eqref{eq:essential}   imply that
		 $$ 
	\lim\limits_{i\to\infty}	\int_{\widetilde{X} } f_{j_i}	\sqrt{-1}^{p^2}	 d\beta \wedge \alpha\wedge \bar{\alpha}\wedge \omega^{n-p-1}=0.
		 $$ 
Recall that $0 \leq f_j \leq 1$,  and $ f_j(x) =1$ for any $x$ satisfying  $d_{\widetilde{X}}(x,x_0)\leq j-1-C$. The above equality together with the monotone convergence theorem imply that 
		\begin{align*} 
		\int_{\widetilde{X} }  	\sqrt{-1}^{p^2}	 d\beta \wedge \alpha\wedge \bar{\alpha}\wedge \omega^{n-p-1}=0.
		\end{align*}  
		It follows \eqref{eq:mutual} that 
\begin{align}\label{eq:pointwise}
\int_{\widetilde{X}}  	\sqrt{-1}^{p^2}	  \alpha\wedge \bar{\alpha} \wedge \omega^{n-p-1}\wedge \pi_X^*f^*T_Y=0. 
\end{align}
%By \Cref{claim:Demailly} again,  $	\sqrt{-1}^{p^2}	\alpha\wedge \bar{\alpha} \wedge \pi_X^*f^*\omega_Y$ is a positive $(p+1,p+1)$-form.  The above equality implies that \begin{align}\label{eq:pointwise}
%\sqrt{-1}^{p^2}	\alpha\wedge \bar{\alpha} \wedge \pi_X^*f^*\omega_Y\equiv 0.
%\end{align} 

\medspace

\noindent {\bf Step 2.}  
We abusively write \( T_Y \) to denote \( \pi_Y^*T_Y \). Let \( y \in \widetilde{Y}^\circ \) be any point, and let \( x \in \widetilde{X} \) be any point in the fiber \( \tilde{f}^{-1}(y) \). Since \( f \) is smooth over \( Y^\circ \), we can find coordinate neighborhoods \( (U; z_1, \ldots, z_n) \) centered at \( x \) and \( (V; w_1, \ldots, w_m) \) centered at \( y \), such that 
\[
\tilde{f}(z_1, \ldots, z_n) = (z_1, \ldots, z_m).
\]
Since \( T_Y|_{Y^\circ} \) is a Kähler form on \( Y^\circ \), after applying a unitary change of coordinates, we may assume that 
\[
\omega(x) = \sqrt{-1} \sum_{j=1}^n dz_j \wedge d\bar{z}_j \quad \text{and} \quad T_Y(y) = \sum_{i=1}^m \lambda_i \sqrt{-1} dw_i \wedge d\bar{w}_i,
\]
where each \( \lambda_i \in \mathbb{R}_{>0} \).
	
	
	For any subset \( I = \{i_1, \ldots, i_p\} \subset \{1, \ldots, n\} \), we denote \( dz_I := dz_{i_1} \wedge \cdots \wedge dz_{i_p} \). Let \( \alpha \in H^{p,0}_{(2)}(\widetilde{X}) \). If we write
	\[
	\alpha|_U = \sum_{|I| = p} \alpha_I \, dz_I, \quad \text{with } \alpha_I \in \mathcal{O}(U),
	\]
	then we have
	\[
	\sqrt{-1}^{p^2} \, \alpha \wedge \bar{\alpha} \wedge \tilde{f}^*T_Y \wedge \omega^{n-p-1}(x)
	= \frac{(n - p - 1)!}{n!}\sum_{|I|=p} \sum_{j \in \{1, \ldots, m\} \setminus I} \lambda_j |\alpha_I(x)|^2 \, \omega(x)^n.
	\]
	By \eqref{eq:pointwise}, it follows that \( \{1, \ldots, m\} \subset I \) for any multi-index \( I \) with \( \alpha_I(x) \neq 0 \). Since \( x \) is arbitrary in the non-empty analytic open subset \( \widetilde{X}^\circ \subset \widetilde{X} \), we conclude that \( \alpha(x) = 0 \) for all \( x \in \widetilde{X}^\circ \) if $p<m$. As \( \alpha \) is a holomorphic \((p,0)\)-form, it follows that \( \alpha \equiv 0 \). This establishes \Cref{item:partial}.
	
	Moreover, if \( n > m \) and \( p \geq m \), then this also proves \eqref{eq:tensor}. Hence, \Cref{item:long} is proved.
	
	
 
 \medspace
 
 \noindent {\bf Step 3.}   
 Let us now prove \Cref{item:d-closed}.  
 For any point \( y \in \widetilde{Y}^\circ \), denote the fiber over \( y \) by \( F_y := \tilde{f}^{-1}(y) \).  
 Every nontrivial \( m \)-form \( \eta \in \Omega^m_{\widetilde{Y}^\circ,y} \) induces a \emph{unique} holomorphic \((p - m)\)-form \( \alpha_y \in H^0(F_y, \Omega^{p-m}_{F_y}) \) such that
 \[
 \alpha = \alpha_y \wedge \tilde{f}^*\eta.
 \]
 This form \( \alpha_y \) can be written explicitly. Using the coordinate charts \( (U; z_1, \ldots, z_{n+m}) \) and \( (V; w_1, \ldots, w_m) \) introduced in Step 2, we may write
 \[
 \eta = \lambda \, dw_1 \wedge \cdots \wedge dw_m,
 \]
 for some \( \lambda \in \mathbb{C}^* \) (note that \( \lambda \) depends on the choice of coordinates).
 
 For any multi-index \( I \supset \{1, \ldots, m\} \), denote by \( \tilde{I} := I \setminus \{1, \ldots, m\} \).  
 Then, by \Cref{item:long}, we have  
\begin{align}\label{eq:express2}
	 \alpha|_U=\sum_{|I|=p,I\supset\{1,\ldots,m\}}\alpha_Idz_{\tilde{I}}\wedge dz_1\wedge \cdots\wedge dz_m.
\end{align} 
Then the holomorphic \((p - m)\)-form \( \alpha_y \) on \( F_y \cap U \) is defined by
\begin{align}\label{eq:express}
	\alpha_y := \lambda^{-1} \cdot \sum_{|I| = p,\; I \supset \{1, \ldots, m\}} \alpha_I|_{F_y} \, dz_{\tilde{I}}.
\end{align}
Since \( \alpha_y \wedge \tilde{f}^*\eta = \alpha \), it follows that the definition of \( \alpha_y \) is independent of the choice of coordinates. Therefore, as we vary the coordinate open subsets, the local forms glue together to define a global holomorphic \((p - m)\)-form \( \alpha_y \) on \( F_y \).

Now, since \( \alpha \) is \( d \)-closed by assumption, the expression for \( \alpha \) implies
\begin{align}\label{eq:together}
	0 = d\alpha = \lambda^{-1} \cdot \sum_{\substack{|I| = p\\ I \supset \{1, \ldots, m\}}} \sum_{j = m + 1}^{n} \frac{\partial \alpha_I}{\partial z_j} \, dz_j \wedge dz_{\tilde{I}} \wedge dz_1 \wedge \cdots \wedge dz_m.
\end{align}
This yields
\[
d\alpha_y = \lambda^{-1} \cdot \sum_{\substack{|I| = p\\ I \supset \{1, \ldots, m\}}} \sum_{j = m + 1}^{n} \left( \frac{\partial \alpha_I}{\partial z_j} \bigg|_{F_y} \right) dz_j \wedge dz_{\tilde{I}} = 0.
\]
This proves \Cref{item:d-closed}, and completes the proof of the theorem. 
	\end{proof}
	
	\iffalse
	\begin{lemma}
		Let $f:X\to Y$ be a proper holomorphic fibration from a complex $n$-dimensional manifold $X$  to a complex manifold $Y$ of dimension $m$ with $n>m$. For some $p>m$, let $\alpha$ be a holomorphic $p$-form on the universal cover $\widetilde{X}$ such that $$\alpha\in H^0(\widetilde{X},\tilde{f}^*\Omega_{\widetilde{Y}}^m\otimes \Omega_{\widetilde{X}}^{p-m}),$$
		where $\tilde{f}:\widetilde{X}\to\widetilde{Y}$ is the lift of $f$ to   universal covers.  Then for any $y_0\in \widetilde{Y}$, there is a coordinate neighborhood $(V;w_1,\ldots,w_m)$ centered at $y_0$ and smooth $(p-m)$-forms $\eta$ on $\tilde{f}^{-1}(V)$  such that
		\begin{thmlist}
			\item the restriction of $\eta$ to the fiber $\tilde{f}^{-1}(y)$  of each $y\in V$ is  a holomorphic $(p-m)$-form;
			\item we have
			$$\alpha=\eta\wedge \tilde{f}^*(dw_1\wedge \cdots\wedge dw_m).$$
		\end{thmlist}
	\end{lemma}
	\begin{proof}
		Since $f:X\to Y$ is proper, we can choose finitely many coordinate open subsets $\{(U_i;z^{(i)}_1,\ldots,z^{(i)}_n)\}_{i=1,\ldots,\ell}$ of $X$ and coordinate neighborhood $(V;w_1,\ldots,w_m)$ centered at $y_0$ such that 
		\begin{enumerate}
			\item each $U_i$ is simply connected;
			\item for each $i\in \{1,\ldots,\ell\}$, we have $f(U_i)\subset V$  and
			$f(z^{(i)}_1,\ldots,z^{(i)}_n)=(z^{(i)}_1,\ldots,z^{(i)}_m)$;
			\item there exists an open neighborhood $W$ of $y_0$ such that $f^{-1}(W)\subset \cup_{i=1,\ldots,\ell}U_i$.    
		\item there exists smooth functions $\{\varrho_i:X\to [0,1]\}_{i=1,\ldots,\ell}$ such that ${\rm supp}(\varrho_i)\subset U_i$, and  the sum $\sum_{i=1}^{\ell}\varrho_i$ has constant value 1 on $f^{-1}(W)$.
		\end{enumerate}
Let us express  
	$\alpha|_{U_i}=\sum_{|I|=p}\alpha^{(i)}_I dz^{(i)}_I$, where $dz^{(i)}_I:=  dz^{(i)}_{i_1}\wedge \cdots \wedge dz^{(i)}_{i_p}$ with  $I$  consisting of $i_1<\ldots<i_p$, 	 
	 For any set $I\supset\{1,\ldots,m\}$,  denote by $\tilde{I}:=I\backslash\{1,\ldots,m\}$.   By our assumption, we have
	\begin{align} 
		\alpha|_{U_i}=\sum_{|I|=p,I\supset\{1,\ldots,m\}}\alpha^{(i)}_Idz^{(i)}_{\tilde{I}}\wedge dz^{(i)}_1\wedge \cdots\wedge dz^{(i)}_m.
	\end{align} 
\begin{claim}
	For any $y\in W$, if $f^{-1}$
\end{claim}
		
		
			\end{proof}
	\fi
	We will need the following consequence of \Cref{thm:vanishing2} in the proof of \Cref{main:kollar}. 
	\begin{corollary}\label{prop:abelian}
		Let \( (X, \omega) \) be a compact Kähler $n$-fold, and let \( f: X \to A \) be a holomorphic map to an abelian variety \( A \) such that \( \dim f(X) = n \). Let \( \widetilde{X}_1 \) be a connected component of the fiber product \( X \times_A \widetilde{A} \), where \( \widetilde{A} \to A \) denotes the universal cover of \( A \). Then, for any infinite Galois cover \( \pi':\widetilde{X}' \to X \) that factors through \( \widetilde{X}_1 \), we have
		\[
		H^{p,0}_{(2)}(\widetilde{X}') = 0 \quad \text{for all } p \in \{0, \ldots, n - 1\}.
		\]
		\end{corollary}
		
		\begin{proof}
	 We choose global linear coordinates \( (z_1, \ldots, z_m) \) on \( \widetilde{A} \) such that the standard Kähler form
		\[
		\omega_{\widetilde{A}} := \sqrt{-1} \sum_{i=1}^{m} dz_i \wedge d\bar{z}_i
		\]
		descends to a Kähler form \( \omega_A \) on \( A \). Since \( \widetilde{X}' \) dominates \( \widetilde{X}_1 \), the map \( f: X \to A \) lifts to a holomorphic map \( g: \widetilde{X}' \to \widetilde{A} \). 
		Set \( g_i := z_i \circ g \) for each \( i \); this defines a holomorphic function on \( \widetilde{X}' \).  
		
		We consider the smooth real \((1,0)\)-form
		\[
		\beta := g^*\left( \sqrt{-1} \, \bar{\partial} \sum_{i=1}^{m} |z_i|^2 \right) 
		= \sqrt{-1} \sum_{i=1}^{m} g_i(x) \, \overline{\partial g_i(x)}.
		\]
		Since each \( dz_i \) descends to a holomorphic 1-form on \( A \), it follows that each \( \partial g_i(x) \) descends to a holomorphic 1-form on \( X \). Therefore, there exists a constant \( C > 0 \) such that
		\[
		|\partial g_i(x)|_{\pi'^* \omega} \leq C
		\quad \text{for all } x \in \widetilde{X}' \text{ and } i = 1, \ldots, m.
		\]
		In particular, the map \( g \) is Lipschitz. Fix a base point \( x_0 \in \widetilde{X}' \). Then there exists a constant \( C' > 0 \) such that
		\[
		|g_i(x)| \leq C' \left( d_{\widetilde{X}'}(x, x_0) + 1 \right)
		\quad \text{for all } x \in \widetilde{X}'.
		\]
		It follows that
		\[
		|\beta(x)|_{\pi'^* \omega} \leq m C C' \left( d_{\widetilde{X}'}(x, x_0) + 1 \right).
		\]
		
		Note that
		\[
		d\beta = (\pi')^* f^* \omega_A,
		\]
		where \( f^* \omega_A \) is a closed semi-positive \((1,1)\)-form on \( X \), which is strictly positive on a non-empty Zariski open subset where \( f \) is immersive.
		
		By \Cref{item:partial}, this implies the desired \( L^2 \)-vanishing result.
	 	\end{proof}
	 
	  \iffalse
	\subsection{Semi-Kähler parabolicity and  fundamental groups} 
	The following fact is obvious. We provide a proof here for the sake of completeness.
	\begin{lemma}\label{lem:equiv}
		Let $(X,\omega_X)$ be a  semi-K\"ahler  parabolic manifold. Then $X$ has   generically large  fundamental group.
	\end{lemma}
	\begin{proof}
	  By \cref{def:wkp}, there exists a  1-form $\beta$ with $L^1_{\rm loc}$ coefficients, a continuous psh function $\psi$ on $X$  and a semi-K\"ahler form $\omega_{\rm sk}$ on $X$ such that 
		$d\beta+\pi_X^*\hess\psi$ is a real semi-positive $(1,1)$-form satisfying 
		$$d\beta+\pi_X^*\hess\psi\geq \pi_X^*\omega_{\rm sk}.$$ Let $\Xi$ be a Zariski closed proper subset of $X$ such that $\omega_{\rm sk}$ is strictly positive outside $\Xi$. 	Let $Z$ be any positive-dimensional analytic subvariety of $X$ not contained in $\Xi$.  We will prove that ${\rm Im}[\pi_1(Z)\to \pi_1(X)]$ is infinite.
		
		Suppose, for sake of contradiction, that ${\rm Im}[\pi_1(Z)\to \pi_1(X)]$ is finite. Then for any connected component $Z'$ of $\pi_X^{-1}(Z)$, $\pi_X|_{Z'}:Z'\to Z$ is finite \'etale. Hence $Z'$ is compact. Let  $\mu:Y\to Z'$ be a desingularization.   Let $\omega_Y$ be any K\"ahler form on $Y$.  Then we have
		$$
		0=	\int_Y  \mu^*(d\beta+\pi_X^*\hess\psi)\wedge\omega_Y^{\dim Y-1}\geq  	\int_Y  \mu^*\pi_X^*\omega_{\rm sk}\wedge\omega_Y^{\dim Y-1}>0,
		$$
		which leads to a contradiction. Hence ${\rm Im}[\pi_1(Z)\to \pi_1(X)]$ is infinite. It implies that $X$ has generically large fundamental group.   
	\end{proof}
	In \cref{thm:main2}, we will prove that for   complex projective manifolds with   linear fundamental groups, the converse of the above lemma is also true. This is the main content of \Cref{sec:linear}.
	\fi
	%\subsection{$d$-sublinear}	
	%	\begin{definition}
		%	A differential form $\alpha$ on a complete non-compact Riemannian manifold is called $d$-sublinear  if there exist a differential form $\beta$ and a number $c>0$ such that $d \beta=\alpha$ and $|\beta(x)| \leq c(1+\rho(x, x_0))$, where $\rho(x, x_0)$ stands for the Riemannian distance between $x$ and a base point $x_0$.
		%	\end{definition}  
	%	The concept of $d$-sublinearity is both natural and flexible. This can be seen from the following results.
	
	
	%	\begin{theorem}[\cite{CX}]\label{thm:cao}
		%	Let $N^{2 n}$ be a complete non-compact Kähler manifold of complex dimension $n$ whose Kähler form is $d$-sublinear. If $k \neq n$, then any $L^2$ harmonic $k$-form on $N^{2 n}$ is identically zero.
		%	\end{theorem} 
	
	%	Extending Gromov's terminology, Cao-Xavier \cite{CX} proposed the following:
	%	\begin{definition}[Kähler-parabolic]\label{def:kp}
		%	A compact Kähler manifold is Kähler-parabolic if the Kähler form on its universal cover is $d$-sublinear.
		%	\end{definition} 
	
	%	Accordingly, their results imply the Singer conjecture if the compact K\"ahler manifold $M^{2 n}$ is Kähler-parabolic.
	
	
	\section{Constructing  1-forms via harmonic maps to Euclidean buildings}\label{sec:1-form}  
	Let $K$ be a non-archimedean local field of characteristic zero and let $G$ be a reductive algebraic group defined over $K$. There exists a Euclidean building associated with $G$, which is  the enlarged Bruhat-Tits building and denoted by $\Delta(G)$.  We refer the readers to \cite{KP23,Rou23} for the definition and properties of Bruhat-Tits buildings. 
	
	
	
	Let $(V,W,\Phi)$ be the root system associated with $\Delta(G)$. It means that $V$ is a real Euclidean space endowed with a Euclidean metric and $W$ is an affine Weyl group acting on $V$. Namely, $W$ is a semidirect product   $T\rtimes W^v$,  where $W^v$ is the \emph{vectorial Weyl group}, which is a finite group generated by reflections on $V$, and $T$ is a translation group of $V$.  Here $\Phi$ is the root system of $V$.  It is a finite set of $V^*\backslash \{0\}$ such that $W^v$ acts on $\Phi$ as a permutation.  Moreover, $\Phi$ generates $V^*$.  From the reflection hyperplanes of $W$ we obtain a decomposition of $V$ into facets. Let $\cH$ be set of  hyperplanes of $V$ defined by $w\in W$.  The maximal facets, called \emph{chambers}, are the open connected components of $V\backslash \cup_{H\in \cH}$. 
	
	
For any apartment $A$ in $ \Delta(G)$,  there exists an isomorphism $i_A:A\to V$, which is  called a chart. For two charts $i_{A_1}:A_1\to V$ and $i_{A_2}:A_2\to V$, if $A_1\cap A_2\neq\varnothing$,   it satisfies the following properties:
	\begin{enumerate}[label=(\alph*)]
		\item $Y:=i_{A_2}(i_{A_1}^{-1}(V))$ is convex.
		\item \label{item:weyl}There is an element $w\in W$ such that $w\circ i_{A_1}|_{A_1\cap A_2}=i_{A_2}|_{A_1\cap A_2}$.
	\end{enumerate} 
	The charts allow us to map facets into $\Delta(G)$ and their images are
	also called facets. The  axioms guarantee that these
	notions are chart independent. 
	\iffalse
   Consider the root system $\Phi=\{\alpha_1,\ldots,\alpha_m\}\subset V^*-\{0\}$.  %We define  affine functions 
%	$\alpha_{A,i}:=\alpha_i\circ i_{A}(x)$ on $A$ for each $i$.  
		By \Cref{item:weyl}, there exists an element $w\in W$ such that  
		$
		\alpha_k\circ i_{A_2}|_{A_1\cap A_2}=  \alpha_k\circ w\circ i_{A_1}|_{A_1\cap A_2} 
		$ 
		for any $k=1,\ldots,m$. 
		Recall that $W^v$ permutes $\Phi$. It follows that there exist $a_1,\ldots,a_m\in \bR$ 
		and  a permutation $\sigma$ of $m$-elements   such that
\begin{align}\label{eq:compatible} 
	 	\alpha_k\circ i_{A_2}|_{A_1\cap A_2}=  \alpha_k\circ w\circ i_{A_1}|_{A_1\cap A_2} = \alpha_{\sigma(k)}\circ i_{A_1}|_{A_1\cap A_2} -a_k
\end{align} 
		for any $k=1,\ldots,m$.  
		
		We define a function $\tilde{d}_{\Delta(G)}(x,y):\Delta(G)\times \Delta(G)\to \bR_{\geq  0}$ as follows. For any $x,y\in \Delta(G)$, we pick any apartment $A$ containing $x$ and $y$ and such an apartment always exists by the property of  Bruhat-Tits buildings.  We define
	\begin{align}\label{eq:distancegeneral} 
\tilde{d}_{\Delta(G)}(x,y)=\sum_{i=1}^{m}|\alpha_i\circ i_A(x)-\alpha_i\circ i_A(y)|^2. 
\end{align}
By \eqref{eq:compatible}, it does not depend on the choice of $A$.
	\begin{lemma}
		\begin{thmlist}
			\item  \label{item:distance}	$\tilde{d}_{\Delta(G)}(x,y)$ is a distance function on $\Delta(G)$.  
			\item \label{item:distancediff}  There exists a positive  constant $\lambda>1$ such that $$ \lambda^{-1}d^2_{\Delta(G)}(x,y)\leq d^2_{\Delta(G)}(x,y)\leq \lambda d^2_{\Delta(G)}(x,y),$$  
			 where $d_{\Delta(G)}(\bullet,\bullet)$ is the canonical distance function of $\Delta(G)$.   
		\end{thmlist}
	 %  In particular, 	$r_G:\Delta(G)\to \bR_{\geq  0}$ is a Lipschitz continous function on $\Delta(G)$. 
	\end{lemma}
	\begin{proof}
		The first claim follows from the same proof of \cite[Proposition 6.5. (a)]{Rou09}.  
Let $\{x_1,\ldots,x_N\}\subset V^*$ be a unitary coordinate for $V$.     Then for any points $P_1,P_2\in A$,  their  distance in $\Delta(G)$  %we have
\begin{align}\label{eq:distance}
	d_{\Delta(G)}(P_1,P_2)=\sqrt{\sum_{i=1}^{N}|x_i\circ i_A(P_1)-x_i\circ i_A(P_2)|^2}.
\end{align}
Then there exists a constant $m\times N$-matrix $M$ of real numbers of rank $N$,  such that
$$
\begin{bmatrix}
	\alpha_1, & \cdots, & \alpha_m
\end{bmatrix}^T =M\cdot  \begin{bmatrix}x_1, &\cdots, &  x_N\end{bmatrix}^T.
$$ 
Hence for any $x,y\in A$, we have
\begin{align*}
	 d^2_{\Delta(G)}(x,y)= \begin{bmatrix}x_{1}\circ i_A(y)-x_{1}\circ i_A(x), &\cdots, & x_{N}\circ i_A(y)-x_{N}\circ i_A(x)\end{bmatrix} M^T\cdot \\
	 M  \begin{bmatrix}x_{1}\circ i_A(y)-x_{1}\circ i_A(x), &\cdots, & x_{N}\circ i_A(y)-x_{N}\circ i_A(x)\end{bmatrix}^T.
\end{align*} 
Note that there exists a constant $\lambda\in \bR_{>0}$  such that 
$
\lambda {\rm Id}_{N}-M^TM
$ and 	 $
M^*M-	\lambda^{-1}  {\rm Id}_{N}
$ 
are both positive-definite real symmetric  matrix.  The lemma  follows from the above equality7. 
	\end{proof}
	 \fi
	
	%These affine functions $\{\alpha_{A,1},\ldots,\alpha_{A,m}\}$ allow us to eliminate the ambiguity of $u$ in the expression of coordinates due to the choice of apartments. 	 
	
	
	
Let $X$ be a compact K\"ahler manifold and 	let $\varrho:\pi_1(X)\to G(K)$ be a Zariski dense representation.   	By the work of Gromov-Schoen \cite{GS92} (see also \cite{BDDM} for the quasi-projective case), there exists a $\varrho$-equivariant harmonic mapping $u:\widetilde{X}\to \Delta(G)$.  Such $u$ is moreover pluriharmonic.  	We denote by $R(u)$ the regular set of harmonic map $u$. That is, for any $x\in R(u)$, there exists  an open subset $\Omega_{x}$ containing $x$ such that $u(\Omega_{x})\subset A$ for some apartment $A$.   
Since $G(K)$ acts transitively on the apartments of $\Delta(G)$,    we know that $R(u)$ is the pullback of an open subset $X'$ of $X$. By \cite{GS92}, $X\backslash X'$  has \emph{Hausdorff codimension} at least two.  

 We fix  an orthogonal coordinates $(x_1,\ldots,x_N)$ for $V$.  Define   smooth real functions on $\Omega_x$ by setting
\begin{align}\label{eq:ui}
 u_{A,i}:=x_i\circ i_A\circ u.
\end{align}   The pluriharmonicity of $u$ implies that $\hess u_{A,i}=0$. We consider a smooth real semi-positive $(1,1)$-form on $\Omega_x$ defined by 
\begin{align}\label{eq:form}
   \sqrt{-1}\sum_{i=1}^{N}\partial u_{A,i}\circ \db u_{A,i}. 
 \end{align}
By \cite[\S 3.3.2]{Eys04},  such real $(1,1)$-form does not depend on the choice of $A$, and is invariant under $\pi_1(X)$-action. Therefore, it descends to a smooth real closed semi-positive $(1,1)$-form on $X'$.  
 It is shown in \cite{Eys04} that it extends to a  positive closed $(1,1)$-current $T_{\varrho}$  on $X$ with continuous potential.  
\begin{definition}[Canonical current]\label{def:canonical}
	The above closed positive $(1,1)$-current 	$T_\varrho$ on $X$ is called the \emph{canonical current} of $\varrho$.   
\end{definition}  
\begin{remark}
	Although the $\varrho$-equivariant harmonic map $u$ may not be unique, the unicity result in \cite{DM24} implies that the  $(1,1)$-form \eqref{eq:form} is independent of the choice of $u$. Since the current $T_\varrho$ is defined as the minimal extension of this $(1,1)$-form, it follows that $T_\varrho$ is uniquely determined and independent of the choice of $u$.
\end{remark}
By \cite{Eys04,CDY22,DY23}, one has  a proper fibration  $s:X\to S_{\varrho}$ associated with $\varrho$,  which is called \emph{Katzarkov-Eyssidieux reduction map}. It has the following properties. 
\begin{proposition}[\cite{Eys04,CDY22}]\label{thm:KE}
	Let $X$ be a smooth projective variety and let $\rho:\pi_1(X)\to G(K)$ be a  Zariski dense representation where $G$ is a reductive group over  a non-archimedean local field $K$.  Then there exists a proper morphism $s_\rho:X\to S_\rho$ (so-called Katzarkov-Eyssidieux reduction map) onto a normal projective variety with connected fibers  such that  for any irreducible closed subvariety $Z\subset X$, the following properties are equivalent: 
	\begin{enumerate}
		\item\label{item:KZ1} $\rho({\rm Im}[\pi_1(Z_{\rm norm})\to \pi_1(X)])$ is bounded;
		\item \label{item:KZ3} $\rho({\rm Im}[\pi_1(Z)\to \pi_1(X)])$ is bounded;
		\item \label{item:KZ2}  $s_\rho(Z)$ is a point.
	\end{enumerate}  
	Moreover, there exists a $(1,1)$-current $T_\varrho'$ with continuous potential on $S_\varrho$ such that $s_\varrho^*T_\varrho'=T_\varrho$. \qed
\end{proposition}

	\begin{proposition}\label{lem:equal}
Let $x_0$ be a fixed base point in $\widetilde{X}$. For the canonical current $T_\varrho$ defined in \cref{def:canonical}, we have
	\begin{align}\label{eq:equal} 
 \hess d^2_{\Delta(G)}(u(x),u(x_0)) \geq\pi_X^*T_\varrho,
	\end{align} 
	where $d_{\Delta(G)}(\bullet,\bullet)$ is the distance function on $\Delta(G)$. 
%Moreover, the above equality holds  over a dense open subset $\widetilde{X}^\circ$ such that $\widetilde{X}\backslash \widetilde{X}^\circ$ has zero Lebesgue measure.  
\end{proposition}
\begin{proof}
We assume that $x_0\in R(u)$.  Let $\omega$ be a K\"ahler metric on $X$. 	For any $x\in R(u)$ where $R(u)$ is the regular set of the harmonic map $u$, there exists  an open subset $\Omega_{x}$ containing $x$ such that $u(\Omega_{x})\subset A$ for some apartment $A$. 
	
Note that $ {d}_{\Delta(G)}$ is   $G(K)$-invariant.  Let $\iota:C\to \Omega_x$ be any holomorphic curve.  Then by \cite[Proposition 2.2 in p. 191]{GS92},  we have
	$$
	\Delta d^2(u\circ \iota (x), u(x_0) )\geq |\nabla u\circ \iota|^2,
	$$
	where the  $|\nabla u\circ \iota|^2$ is the norm defined with respect to  $d_{\Delta(G)}$ and $\omega$.  It follows that 
	$$
	|\nabla u\circ \iota|^2 \iota^*\omega= \iota^*(\sqrt{-1}\sum_{i=1}^{N}\partial u_{A,i}\wedge\db u_{A,i})
	$$
	where $u_{A,i}$ is defined in \eqref{eq:ui}. 
	Therefore, over $\Omega_x$ we have
	$$\hess  {d}_{\Delta(G)}^2(u(x), u(x_0) )\geq \sqrt{-1}\sum_{i=1}^{m}\partial u_{A,i}\wedge\db u_{A,i}.$$
By \cref{def:canonical}, we have  \eqref{eq:equal} over the whole $R(u)$. 
	\begin{claim}
\eqref{eq:equal} holds over the whole $\widetilde{X}$.
	\end{claim}
	\begin{proof}
		 Since $u$ is  Lipschitz, it follows that   $d^2_{\Delta(G)}(u(x),u(x_0))$ is a continuous function on $\widetilde{X}$. Recall that $T_\varrho$ is a positive closed $(1,1)$-current with continuous potential. This implies that for any point $x\in \widetilde{X}$, there exist a neighborhood $\Omega_x$ and a continuous function $\phi$ on $\Omega_x$ such that 
		 $$
\hess d^2_{\Delta(G)}(u(x),u(x_0))- \pi_X^*T_\varrho=\hess \phi 
		 $$
		 on $\Omega_x$. Note that $\hess\phi\geq 0$ over $\Omega_x\cap R(u)$. Since the complement of $\Omega_x\cap R(u)$ has Hausdorff codimension at least two in $\Omega_x$, we apply the extension theorem in \cite[Theorem 3.1(i)]{Shi72} to conclude that there is a  psh function $\phi'$ on $\Omega_x$ such that $\phi'|_{\Omega_x\cap R(u)}=\phi|_{\Omega_x\cap R(u)}$.    Since  $\phi$ is continuous,  and $\phi'$ is upper semi-continuous, it follows that   $\phi=\phi'$ on $\Omega_x$.  This shows that \eqref{eq:equal} holds over $\Omega_x$, hence over the whole $\widetilde{X}$.  The claim is proved. 
	\end{proof}
		The proposition is proved.   
	%  Let $\widetilde{X}^\circ$ be the set of points $x$ in $\widetilde{X}$ such that there exists an open neighborhood $\Omega_x$ containing $x$ such that $u(\Omega_x)$ is contained in the closure of a chamber $\cC$ of the building $\Delta(G)$.  By \cite[Proposition 4.15]{DM24}, $\widetilde{X}^\circ$ is a dense open subset such that $\widetilde{X}\backslash \widetilde{X}^\circ$  has zero Lebesgue measure. Note that there exists an apartment $A$ containing both $\overline{\cC}$ and $x_0$. It follows that for any $y\in \Omega_{x}$, we have
%\begin{align}\label{eq:distanced}
%	d^2_{\Delta(G)}(u(y),u(x_0))=\sum_{i=1}^{N}|u_{A,i}(y)-u_{A,i}(x_0)|^2.
%\end{align}
%Therefore, 
%	$$
%	\hess	d^2_{\Delta(G)}(u(y),u(x_0))=  \sqrt{-1}\sum_{i=1}^{m}\partial u_{A,i}(y)\wedge  \db u_{A,i}(y)
%	$$
%	by the pluriharmonicity of $u_{A,i}$.  
%	This proves that over $\widetilde{X}^\circ$,	 we have
%	\begin{align*}  
%		\hess d^2_{\Delta(G)}(u(x),u(x_0)) = \pi_X^*T_\varrho(x).
%	\end{align*}   
\end{proof}	  
%\begin{remark}
%	Note that we cannot expect the equality \eqref{eq:equal} holds over $R(u)$.  Here is an example.  	Consider  a tree $T\subset  \bR^2$ defined by $T=\bR\times \{0\}\cup 0\times \bR_{\geq 0}$ and thus $(0,0)$ is the vertice of $T$.  Consider a pluriharmonic  map   $u:\mathbb{D}\to   T$ defined by
%	 $
%	z\mapsto  ({\rm Re}(z),0).
%	$ 
%	Let  $P:=(0,1)$ be a base point in $T$. Then $d^2(u(z),P)=(|{\rm Re}(z)|+1)^2$, where $d(\bullet,\bullet)$ denotes the distance function on the tree $T$.  In this case, we note that $\hess d^2(u(z),P)$ does not has absolutely continuous coefficients on $\mathbb{D}$: on the line $\mathbb{D}\cap ({\rm Re}(z)=0)$, the coefficients of  $\hess d^2(u(z),P)$ have non-trivial singular part for its Lebesgue decomposition. However, the regular locus  $R(u)$ is the whole disk $\mathbb{D}$. 
%\end{remark}


Let us denote by
\begin{align}\label{eq:defbeta}
	\beta_\varrho(x):=\sqrt{-1}\db d^2_{\Delta(G)}(u(x),u(x_0)).
\end{align}
Since $d^2_{\Delta(G)}(u(x),u(x_0))$ is a psh function on $\widetilde{X}$ by \cref{lem:equal}, by \cite[Theorem 1.46]{GZ17}, $\beta_\varrho$ has $L_{\rm loc}^1$-coefficients. 
 \begin{lemma}\label{lem:distance}
 There exists a dense open subset $\widetilde{X}^\circ $ of $\widetilde{X}$ whose complement has zero Lebesgue measure such that $\beta_\varrho$ is smooth. Moreover, there exists  	a number $c>0$ such that for any $x\in \widetilde{X}^\circ$, we have
 	\begin{align*}
 		|\beta_\varrho(x)|_{\pi_X^*\omega} \leq  c(1+d_{\widetilde{X}}(x, x_0)).
 	\end{align*}  Here $\omega$ is a K\"ahler metric on $X$.
 \end{lemma}
 \begin{proof}
 Let $\widetilde{X}^\circ$ be the set of points $x$ in $\widetilde{X}$ such that there exists an open neighborhood $\Omega_x$ containing $x$ such that $u(\Omega_x)$ is contained in the closure of a chamber $\cC$ of the building $\Delta(G)$.  By \cite[Proposition 4.15]{DM24}, $\widetilde{X}^\circ$ is a dense open subset such that $\widetilde{X}\backslash \widetilde{X}^\circ$  has zero Lebesgue measure. Note that there exists an apartment $A$ containing both $\overline{\cC}$ and $x_0$. It follows that for any $y\in \Omega_{x}$, we have
 \begin{align}\label{eq:distanced}
 	d^2_{\Delta(G)}(u(y),u(x_0))=\sum_{i=1}^{N}|u_{A,i}(y)-u_{A,i}(x_0)|^2.
 \end{align}
% For  any $x\in \widetilde{X}^\circ$, it has an open neighborhood $\Omega_x$ and    an apartment $A$ such that $u(\Omega_x)\subset A$ and $u(x_0)\in A$.  %, and $f_{A,i}=x_i\circ u$, 
 Then  by \eqref{eq:distanced} and \eqref{eq:defbeta}, one has
 	$$
 \beta_\varrho(y)= \sqrt{-1} \sum_{i=1}^{N} (u_{A,i}(y)-u_{A,i}(x_0))\db u_{A,i}(y).
 	$$
 	Since $u$ is   Lipschitz and $\varrho$-equivariant,  it follows that there exists a uniform  constant $c_1>0$ such that for any $y\in \widetilde{X}$, we have $$d_{\Delta(G)}(u(y),u(x_0))\leq  c_1(1+d_{\widetilde{X}}(y, x_0)).$$  
 Note that for any $y\in \Omega_x$, we have
 	$$
 	|u_{A,i}(y)-u_{A,i}(x_0)|\leq  d_{\Delta(G)}(u(y),u(x_0)).
 	$$ 
 	On the other hand, by the Lipschitz condition of $u$, there exists another uniform constant $c_2>0$ such that for any $y\in \Omega_x$, we have
 	$$
 	\sum_{i=1}^{N}|\db u_{A,i}(y)|_{\pi^*\omega}\leq  c_2.
 	$$
 	In conclusion, we have
 	\begin{align*}
 		|\beta_\varrho(y)|_{\pi_X^*\omega} \leq    c_1c_2(1+d_{\widetilde{X}}(y, x_0))\quad \mbox{for any }y\in \Omega_x.
 	\end{align*} 
 	Since $x$ is any point in $\widetilde{X}^\circ$, the above inequality holds for any $x\in \widetilde{X}^\circ$.  The lemma is proved. 
 \end{proof}  
	
	%%%%%%%%%%
	\iffalse
	\subsection{On the Carlson-Toledo conjecture}
	A famous conjecture by Carlson-Toledo predicts the following: if $\Gamma$ is an infinite Kähler group, then 
	$H^2(\Gamma,\bR)\neq 0$. It has the following equivalent interpretation (see \cite{KKM11}). 
	\begin{conj}[Carlson-Toledo]\label{conj:CT}
		Let $X$ be a compact Kähler manifold with infinite fundamental group. There exists a finite étale cover $X^{\prime}$ of $X$ such that the natural morphism $H_2(\widetilde{X}, \mathbb{R}) \longrightarrow H_2\left(X^{\prime}, \mathbb{R}\right)$ is not surjective, or, equivalently, the natural morphism $H^2\left(X^{\prime}, \mathbb{R}\right) \longrightarrow H^2(\widetilde{X}, \mathbb{R})$ is not injective.
	\end{conj} 
	Reznikov and Klingler-Koziarz-Maubon proved the following result: 
	\begin{theorem}[\protecting{\cite[Corollary 1.4]{KKM11}}]\label{KKM}
		Let $\Gamma$ be an infinite Kähler group.  If $\Gamma$ does not satisfy \Cref{conj:CT}, then for any representation $\varrho:\pi_1(X)\to \GL_N(\bC)$ one has $H^1(\Gamma,\mathrm{ad}\varrho)=0$.  
	\end{theorem} 
	As a byproduct of  \cref{lem:equal}, we can  prove the following:
	\begin{theorem}
		Let $\Gamma$ be an infinite K\"ahler group. 
		%  Let $X$ be a compact Kähler manifold with infinite fundamental group.  
		If $\Gamma$ does not satisfy \Cref{conj:CT}, then any semisimple representation $\varrho:\pi_1(X)\to \GL_N(\bC)$ is rigid and integral.
	\end{theorem}
	\begin{proof}
		Assume that $X$ is a compact Kähler manifold, and $\Gamma=\pi_1(X)$. We first prove that the Betti moduli space $M_{\rm B}(X,N)$ (cf. \cite{Sim94}) is zero dimensional for any $N$. Otherwise, if $M_{\rm B}(X,N)$ is positive dimensional, then by \cite{Yam10}, there exists a non-archimedean local field of characteristic 0, and a linear and unbounded representation $\tau:\pi_1(X)\to \GL_N(K)$. By \cite{DY23}, the semisimplification $\sigma:\pi_1(X)\to \GL_N(\bar{K})$ of $\tau$ is also unbounded.  Let $L$ be a finite extension such that $\sigma:\pi_1(X)\to \GL_N(L)$. The canonical current $T_\sigma$ is non-trivial since $\sigma$ is unbounded.   Note that $T_\sigma$ is smooth at a general point of $X$. Hence $\int_{X}T_\sigma\wedge \omega^{n-1}>0$ where $n=\dim X$.  It follows that its cohomology class $\{T_{\sigma}\}\in H^2(X,\bR)$ is non-trivial.  By \cref{lem:equal}, $\pi_X^*T_\sigma$ is $d$-exact. It follows that $\pi_X^*\{T_\sigma\}$ is zero in $H^2(\widetilde{X},\bR)$.  \Cref{conj:CT} is satisfied. Hence $M_{\rm B}(\pi_1(X),N)$ is zero dimensional. 
		
		Therefore, any $\varrho:\pi_1(X)\to \GL_N(\bC)$ is rigid. It is thus conjugate to some representation $\varrho':\pi_1(X)\to \GL_N(K)$, where $K$ is a number field. Let $\nu$ be a non-archimedean place of $K$ and let $K_\nu$ be the non-archimedean completion of $K$ with respect to $\nu$. If $\varrho':\pi_1(X)\to \GL_N(K_\nu)$ is unbounded, then the above arguments shows that the canonical current $T_{\varrho'}$ of $\varrho'$ is non-trivial in $H^2(X,\bR)$, but its lift  $\pi_X^*T_{\varrho'}$ is trivial in $H^2(\widetilde{X},\bR)$. \Cref{conj:CT} is also satisfied. Hence $\varrho':\pi_1(X)\to \GL_N(K_\nu)$ is  bounded  for each non-archimedean place  $\nu$.  Hence $\varrho':\pi_1(X)\to \GL_N(\cO_K)$, which is integral.   The theorem is proved.  
	\end{proof}
	\fi
	
	\section{$L^2$-vanishing theorem and generically large  local systems}\label{sec:linear}
 In this section we will prove \cref{main:kollar}.  In \Cref{subsec:semisimple},   we address the case where $\varrho$ is semisimple, utilizing techniques from the proof of the reductive Shafarevich conjecture in   \cite{Eys04,DY23}.  The desired 1-form $\beta$, as required in \Cref{thm:vanishing2},  arises from   1-forms $\beta_\tau$ associated with $\tau:\pi_1(X)\to \GL_N(K)$ defined in \eqref{eq:defbeta}, where $K$ is a non-archimedean local field.   We then reduce the proof   to \Cref{item:partial}.
	
	
 For the proof of the general cases of \cref{main:kollar}, we will apply techniques from the proof of the linear Shafarevich conjecture in \cite{EKPR12}.   
 Using similar techniques as in the semi-simple case,   we first  construct a suitable fibration $f:X\to Y$ (the reductive Shafarevich morphism ${\rm sh}_M^0:X\to {\rm Sh}_M^0(X)$)  such that the conditions in \cref{thm:vanishing2} are fulfilled. Moreover, by the structure of the linear Shafarevich morphism, for  \emph{almost every}  smooth fiber  of $f$, the conditions in \cref{prop:abelian} are satisfied.  
 
 Therefore, if there exists a non-trivial $\alpha\in H^{p,0}_{(2)}(\widetilde{X})$ for some $p\in \{0,\ldots,\dim X-1\}$,  we can  apply \cref{item:partial}  to show that $p\geq m:=\dim Y$, and use \cref{item:d-closed} to obtain a non-trivial $L^2$ holomorphic $(p-m)$-form on the universal cover of  at least one  smooth fiber  of $f$. Finally, we apply \cref{prop:abelian} to obtain a contradiction.  % contradiction.  Let us assume by contradiction that there exists non-trivial $\alpha\in H^{(p,0)}(\widetilde{X})$ for some $0\leq p<\dim X$.  % Inspired by the concept of the \emph{Betti form} on the universal family of Abelian varieties (cf. \cite{CGHX21}), we will use some type of ``normal functions'' on $\widetilde{X}$ induced by the $\bR$-VMHS $\sM$  to construct a smooth $(1,1)$-form $\omega_\sM$  on  $X$.   However, such $\omega_\sM$ is not  semi-positive in general. We   show that it is nevertheless semi-positive along the fiber direction of the \emph{reductive Shafarevich morphism}. Combining such vertical positivity of $\omega_\sM$ and the horizontal positivity constructed in the proof of \cref{thm:main}, we will obtain the desired positivity.  The more delicate issue is the construction of the 1-form $\beta$ and exhausting smooth function $r$ on the universal cover $\widetilde{X}$ satisfying the conditions in \cref{def:wkp}.   This is the main content of \Cref{subsec:mono}. 
	% will then utilize the positivity derived from the semisimplification of $\varrho$ to compensate for the possible horizontal negativity of $f^*\omega_\cM$ in the general case where $f$ is not contained within the fiber of $\cM \to \cD$. 
	
 	\Cref{subsec:semisimple} is covered by \Cref{subsec:linear}, and readers can skip it if they prefer to proceed directly to the proof of the general case of \cref{main:kollar}. 
 	\subsection{Two technical results}
 	 	Note that in \eqref{eq:defbeta}, we construct certain  1-forms $\beta_{\tau_i}$ on $\widetilde{X}$ associated with the non-archimedean representations $\tau_i$   in \cref{thm:Sha}. For the $\bC$-VHS $\cL$, there is also another way to construct  similar 1-forms. This was established by Eyssidieux in \cite{Esy97} and we recollect some facts  therein. 
 	\begin{proposition}[\protecting{\cite[Proposition 4.5.1]{Esy97}}]\label{prop:eys}
 		Let $X$ be a smooth projective variety and let $\cL$ be a $\bC$-VHS on $X$. Let $\cD$ be the period domain of $\cL$ and let $f:\widetilde{X}\to \cD$ be the period map. Let $q:\cD\to \cR$ be the natural quotient where $\cR$ is the corresponding Riemannian symmetric space of $\cD$. Define $\omega_\cL:=\sqrt{-1}{\rm tr}(\theta\wedge\theta^*)$ as in \cref{thm:Sha}, which is a smooth closed positive $(1,1)$-form. Then there exists a smooth  function $\psi_{\cD}:\cR\to \bR_{>0}$ and  constants $c>0$   depending  only on $\cD$  such that  the function $\phi:=\psi_{\cD}\circ q\circ f$ is smooth and plurisubharmonic and we have
 		\begin{align}\label{eq:eys}
 			\sqrt{-1}\partial\phi\wedge\db\phi&\leq \pi_X^*\omega_\cL,\\\label{eq:eys2}
 			c\pi_X^*\omega_\cL&\leq 	\hess \phi.
 		\end{align}  \qed
 	\end{proposition} 
 	Therefore, we have the following consequence.
 	\begin{corollary}\label{cor:eta}
	Let $X$ be a smooth projective variety and let $\cL$ be a $\bC$-VHS on $X$. Fix a Kähler metric $\omega$ on $X$. For  $\bC$-VHS $\cL$ on $X$, there exists a smooth 1-form $\eta_\cL$ on the universal cover $\widetilde{X}$ and a constant $C > 0$ such that, $d\eta_\cL$ is a real positive $(1,1)$-form satisfying
 \begin{align}
	|\eta_\cL|_{\pi_X^*\omega} \leq C, \quad d\eta_\cL \geq \pi_X^*\omega_\cL.
\end{align}  
\end{corollary}
\begin{proof}
	Define
	\[
	\eta_\cL := \frac{\sqrt{-1}  \db \phi}{c},
	\]
	where $\phi$ and the constant $c > 0$ are as in \cref{prop:eys}. Since there exists a constant $c_1 > 0$ such that $\omega \geq c_1 \omega_\cL$, it follows from \eqref{eq:eys} that
	\[
	|\eta_\cL|_{\pi_X^*\omega} \leq \frac{c_1}{c}.
	\]
	Moreover, by \eqref{eq:eys2}, we have
	\[
	d\eta_\cL \geq \pi_X^*\omega_\cL.
	\]
\end{proof}
	\subsection{Case of semi-simple local systems}\label{subsec:semisimple} 
		We will use   techniques in proving the reductive Shafarevich conjecture in \cite{DY23,Eys04}.  We summarize the main results needed in  proving \cref{thm:main3} as follows. 
	\begin{theorem}[\protecting{\cite[Proof of Theorem 4.31]{DY23}}]\label{thm:Sha}
		Let \( X \) be a smooth projective variety, and let \( \varrho: \pi_1(X) \to \GL_N(\bC) \) be a semisimple representation. If \( \varrho \) is generically large, then there exist:
		\begin{enumerate}
			\item a family of Zariski dense representations \( \{\tau_i: \pi_1(X) \to G_i(K_i)\}_{i=1,\ldots,\ell} \), where each \( G_i \) is a reductive group over a non-archimedean local field \( K_i \) of characteristic zero;
			\item a \( \bC \)-VHS \( \cL \) on \( X \),
		\end{enumerate}
		such that
		\[
		T_{\tau_1} + \cdots + T_{\tau_\ell} + \sqrt{-1} \, \mathrm{tr}(\theta \wedge \theta^*)
		\]
		is a closed positive \( (1,1) \)-current on \( X \), which is smooth and strictly positive over a non-empty analytic open subset \( X^\circ \subset X \).
		
		Here:
		\begin{itemize}
			\item \( T_{\tau_i} \) denotes the canonical current on \( X \) associated with \( \tau_i \), as defined in \cref{def:canonical};
			\item \( \theta \) is the Higgs field of the Hodge bundle associated with \( \cL \), and \( \theta^* \) is its adjoint with respect to the Hodge metric. \qed
		\end{itemize}
	\end{theorem}
	\begin{theorem}\label{thm:main3}
		Let $X$ be a smooth projective $n$-fold. Let $\varrho:\pi_1(X)\to \GL_N(\bC)$ be a semisimple representation. 
  If $\varrho$ is  generically large,  then  
 $H^{p,0}_{(2)}(\widetilde{X})=0$   for  $0\leq p\leq n-1$.
	\end{theorem} 
	\begin{proof} 
By \cref{thm:Sha}, there exist	Zariski dense   representations $\{\tau_i:\pi_1(X)\to G_i(K_i)\}_{i=1,\ldots,\ell}$  where each $G_i$  is a reductive algebraic group over a non-archimedean field $K_i$, along with a $\bC$-VHS $\cL$  on $X$ satisfying the stated properties.   
We denote by $\omega_\cL:=\sqrt{-1}{\rm tr}(\theta\wedge\theta^*)$, that  is a smooth $(1,1)$-form defined in \cref{thm:Sha}.  We fix a K\"ahler form $\omega_X$ on $X$, and, by abuse of notation, also denote by $\omega_X$ its pullback to the universal cover $\widetilde{X}$.
	 % and we have
		%	$$
		%	\mu^*\omega'+\hess \mu^*\psi=i\sum_{i=1}^{m}\eta_i\wedge\overline{\eta_i}+\mu^*\omega_\cL.
		%	$$
		%	Therefore, $\mu^*(i\db \psi)$ is a smooth 1-form on $\xsp$, where \emph{a priori} $i\db \psi$ is defined as a distribution.  Let $\omega_2$ be a smooth K\"ahler metric on $\xsp$. Then there exists a constant $c_1>0$ such that $\mu^*\omega\leq c_1\omega_2$.   It follows that  outside the ramification locus of $\mu$, we have
		%\begin{align}\label{eq:es1}
		% 	\mu^* | i\db \psi|_{\omega} \leq   |\mu^*i\db \psi|_{\omega_2}\leq c_1
		%\end{align} 
		%	where $c_1$ is a positive constant.  In particular, $\db\psi$ is a $(0,1)$-form with local coefficients being  bounded functions. 
		
	 In what follows, for any form $\eta$ we shall denote by $|\eta|$ its norm with respect to $\omega_X$.   		Let $u_i:\widetilde{X}\to \Delta(G_i)$ be the $\tau_i$-equivariant pluriharmonic map whose existence is ensured by \cite{GS92}, where $\Delta(G_i)$ is the Bruhat-Tits building of $G_i$.  By \eqref{eq:defbeta},  if we define
		\begin{align}\label{eq:beta}
		\beta_i(x)=\sqrt{-1} \db d^2_{\Delta(G_i)}(u_i(x),u_i(x_0)),
		\end{align}   then  by \cref{lem:distance}, $\beta_i$ has $L_{\rm loc}^1$-coefficients and  is  smooth outside a set of zero Lebesgue measure.  Moreover,     there exists a constant $c_1>0$  with 
		$
	|\beta_i(x)|\leq_{\rm a.e.}c_1 (d_{\widetilde{X}}(x,x_0)+1) 
		$  
		for any $i$. 
	
		
		Let $\eta_\cL$ be  smooth 1-form in $\widetilde{X}$ defined in \cref{cor:eta}. Then we have 
		\begin{align} \label{eq:between}
|\eta_\cL|\leq c_2,\quad d\eta_\cL\geq \pi_X^*\omega_\cL 
		\end{align}  
		for some positive constant $ c_2$.    
	Define 
 $\beta:=\eta_\cL+\sum_{i=1}^{\ell}\beta_i$.  Then $\beta$ has $L_{\rm loc}^1$-coefficients. By \cref{lem:distance} and \eqref{eq:between}, it satisfies 
\begin{align}\label{eq:dsub}
	  |\beta(x)|\leq_{\rm a.e.} c_3(1+d_{\widetilde{X}}(x, x_0)) 
\end{align} 
	 for some constant $c_3>0$.   	By \cref{lem:equal} and \eqref{eq:between}, 	  we have
	 \begin{align*} d\beta\geq   \pi_X^*(\sum_{i=1}^{\ell}T_{\tau_i}+\omega_{\cL} ). 
	 \end{align*}   
 By  \cref{thm:Sha}, $\sum_{i=1}^{\ell}T_{\tau_i}+\omega_{\cL}$ is closed positive $(1,1)$-current, that is smooth and strictly positive over an analytic   open subset of $X$.  Therefore,   the 1-form $\beta$ on $\widetilde{X}$ satisfies the assumptions of \cref{thm:vanishing2}, with $f$ taken to be the identity map. We thus obtain the desired $L^2$ vanishing result.
	\end{proof} 
	\begin{remark}\label{rem:final}
		If we compare the proof of \cref{thm:main3}  with that  of the reductive Shafarevich conjecture in \cite{DY23,Eys04}, we observe striking similarities in their approaches. 
		Indeed, in \cite[Proposition 4.1.1]{Eys04}, Eyssidieux proved the following result: let $X$ be   a compact K\"ahler normal variety. If there exist a continuous plurisubharmonic function $\phi:\widetilde{X}\to \bR_{>0}$ and a positive closed $(1,1)$-current $T$ on $X$ with continuous potential such that  $\{T\}$ is a K\"ahler class and  $\hess\phi\geq \pi_X^*T$, then $\widetilde{X}$ is Stein. Therefore, if $X$ is a smooth projective variety endowed with a  semisimple and large representation $\varrho:\pi_1(X)\to \GL_N(\bC)$, then we can prove that $\widetilde{X}$  is Stein using \cref{thm:Sha} as follows. 
		
		Consider the continuous function
		$$
		\phi_0:=\sum_{i=1}^{\ell}	 d^2_{\Delta(G_i)}(u_i(x),u_i(x_0))+\frac{\phi}{c_2}
		$$
		on $\widetilde{X}$, where $u_i$ and $\phi$ are defined  in the proof of \cref{thm:Sha}. We then have
		$$
		\hess\phi_0\geq \pi_X^*(\sum_{i=1}^{\ell}T_\ell+ \omega_\cL).
		$$
		Note that $\{ \sum_{i=1}^{\ell} T_\ell+\omega_\cL\}$ is a K\"ahler class in $X$ if $\varrho$ is large by \cref{thm:Sha}. Therefore, by the above Eyssidieux's criterion, we conclude that $\widetilde{X}$ is Stein. 	
 %		However, as we shall see in the rest of the paper, the proof of \Cref{main:kollar} for the case where $\varrho$ is linear is quite involved. We need to develop some new methods to apply the techniques used in the proof of the linear Shafarevich conjecture.
	\end{remark}
	
	
	 
 
	
	 


\subsection{Proof of \cref{main:kollar}}\label{subsec:linear}
In this subsection we will prove \cref{main:kollar}. 
\begin{theorem} \label{thm:main2}
	Let $X$ be a smooth projective $n$-fold. Let $\varrho:\pi_1(X)\to \GL_N(\bC)$ be a generically large  representation.  Then
  $H^{p,0}_{(2)}(\widetilde{X})=0$   for  $0\leq p\leq n-1$.
\end{theorem}
\begin{proof}
	We fix a smooth K\"ahler metric $\omega_X$ on $X$. To lighten the notation, we write  $\omega_X$   instead of   $\pi_X^*\omega_X$ abusively.  
	
	
	\noindent {\bf Step 1.} Consider the Betti moduli space \( M := M_{\mathrm{B}}(X, \GL_N)\). Let 
$
	\mathrm{sh}_M : X \to \mathrm{Sh}_M(X)
$ 
	denote the reductive Shafarevich morphism associated with \( M \) (cf.\ \cite{Eys04, DY23}). In other words, it is a morphism over a normal projective variety such that, for any closed subvariety \( Z \subset X \),  \( \mathrm{sh}_M(Z) \) is a point if and only if, for every \emph{semisimple} representation \( \sigma: \pi_1(X) \to \mathrm{GL}_N(\mathbb{C}) \), the image \( \sigma(\mathrm{Im}[\pi_1(Z) \to \pi_1(X)]) \) is finite.
	
	By \cite[Proof of Theorem 3.29]{DY23}, after replacing \( X \) by a finite étale cover, there exist:
	\begin{itemize}
		\item a family of Zariski dense representations \( \{\tau_i : \pi_1(X) \to G_i(K_i)\}_{i=1,\ldots, \ell} \), where each \( G_i \) is a reductive group over a non-archimedean local field \( K_i \) of characteristic zero;
		\item a complex variation of Hodge structure \( \cL \);
		\item a smooth semi-positive \( (1,1) \)-form \( \omega_M \) on \( \mathrm{Sh}_M(X) \);
		\item for each \( i \), a closed positive \( (1,1) \)-current \( S_i \) on \( \mathrm{Sh}_M(X) \).
	\end{itemize}
	These data satisfy
	\[
	T_{\tau_i} = \mathrm{sh}_M^* S_i, \quad \text{and} \quad \sqrt{-1} \, \mathrm{tr}(\theta \wedge \theta^*) = \mathrm{sh}_M^* \omega_M,
	\]
	where \( T_{\tau_i} \) denotes the canonical current associated with \( \tau_i \) in \cref{def:canonical}. 

Moreover, there exists a non-empty analytic open subset \( U \subset \mathrm{Sh}_M(X) \) such that the positive \( (1,1) \)-current
$
\sum_{i=1}^\ell S_i + \omega_M
$ 
is smooth and strictly positive on \( U \). We denote by \( f : X \to Y \) the reductive Shafarevich morphism \( \mathrm{sh}_M \). 
	
	\medspace
	
		\noindent {\bf Step 2.}  
	For each $i$, let  $\beta_i$ be  the 1-form on $\widetilde{X}$ defined in \eqref{eq:beta} associated with the representation $\tau_i$. By \cref{lem:equal}, it is smooth outside a set of Lebesgue measure zero and satisfies 
	\begin{align}\label{eq:f1}
		d\beta_i\geq \pi_X^*T_{\tau_i}.
	\end{align}	 	By \cref{lem:distance}, 	 there exists  	a number $c_0>0$    such that for any $i\in\{1,\ldots,\ell\}$, we have
	\begin{align}\label{eq:sublinear4}
		|\beta_i(x)|_{\omega_X} \leq_{\rm a.e.} c_0(1+d_{\widetilde{X}}(x, x_0)).
	\end{align}   
	
	Let $\eta_\cL$ be  smooth 1-form in $\widetilde{X}$ defined in \cref{cor:eta}. Then we have 
	\begin{align} \label{eq:estimate5}
		|\eta_\cL|\leq c_1,\quad d\eta_\cL\geq \pi_X^*\sqrt{-1}\mathrm{tr}(\theta \wedge \theta^*) 
	\end{align}  
	for some positive constant $ c_1$.    
	Write 
	$$\beta:=\beta_1+\cdots+\beta_\ell+\eta_\cL. $$ Then $\beta$ has $L_{\rm loc}^1$-coefficients and smooth  outside a set of Lebesgue measure zero.  By \eqref{eq:sublinear4} and \eqref{eq:estimate5}, we have
	\begin{align}\label{eq:bound}
		|\beta(x)|_{\omega_X}\leq_{\rm a.e.} c_2(1+d_{\widetilde{X}}(x, x_0)),
	\end{align} 
	for some constant $c_2>0$. 
	By  \eqref{eq:f1} and \eqref{eq:estimate5}, one has \begin{align}\label{eq:kahlerpara}
		\pi_X^*f^*( \sum_{i=1}^\ell S_i + \omega_M)=	T_{\tau_1}+\cdots+T_{\tau_\ell}+\sqrt{-1}{\rm tr}(\theta\wedge\theta^*)\leq  	  d\beta.   
	\end{align}  
	Recall that there exists a non-empty analytic open subset \( U \subset Y \) such that the positive \( (1,1) \)-current
	$
	\sum_{i=1}^\ell S_i + \omega_M
	$
	is smooth and strictly positive on \( U \).
	Therefore,   the 1-form $\beta$ satisfies   the conditions in \cref{thm:vanishing2}.  
	
	\medspace
	
	\noindent {\bf Step 3.} Let \( H \) be the intersection of the kernels of all representations \( \pi_1(X) \to \GL_N(\bC) \). By \cite[Theorem~1]{EKPR12}, the Shafarevich morphism 
	\[
	\mathrm{sh}_H: X \to \mathrm{Sh}_H(X)
	\]
	associated with \( H \) exists. That is, \( \mathrm{sh}_H \) is a morphism to a normal projective variety with connected fibers, characterized by the following property: for any subvariety \( Z \subset X \), the image 
	\[
	\mathrm{Im}[\pi_1(Z) \to \pi_1(X)/H]
	\]
	is finite if and only if \( \mathrm{sh}_H(Z) \) is a point.
	
	In particular, since the representation \( \varrho: \pi_1(X) \to \GL_N(\bC) \) is generically large, it follows that \( \mathrm{sh}_H \) is a birational morphism. Moreover, the morphism \( \mathrm{sh}_M \) factors through \( \mathrm{sh}_H \)
	\begin{equation*}
		\begin{tikzcd}
		X\arrow[r,"{\rm sh}_H"] \arrow[dr, "{\rm sh}_M"' ]& \mathrm{Sh}_H(X)\arrow[d,"q"]\\
&		 \mathrm{Sh}_M(X)
		\end{tikzcd}
	\end{equation*}
	
	By \cite[\S5]{EKPR12}, for a general smooth fiber \( Z \) of \(  \mathrm{sh}_H \), there exists an abelian variety \( A \) and a morphism 
	\[
	a: Z \to A
	\]
	such that the Stein factorization of \( a \) coincides with the restriction \( \mathrm{sh}_H|_Z: Z \to \mathrm{sh}_H(Z) \). Since \( \mathrm{sh}_H \) is birational, we conclude that
	\[
	\dim Z = \dim a(Z).
	\]
	
	\medspace
	

	
	\medspace
	
	\noindent {\bf Step 4.} We will apply \cref{thm:vanishing2} and use its notations as defined therein, without re-explaining their meanings.  Assume by contradiction that, for some $p\in \{0,\ldots,n-1\}$, there is a non-trivial $\alpha\in H^{p,0}_{(2)}(\widetilde{X})$. By \cref{item:partial},  we have $n>m$ and $p\geq m$. Furthermore, over the analytic open subset $Y^\circ$ of  $Y^{\rm reg}$, 
	we have   \begin{align*} 
		\alpha|_{\widetilde{X}^\circ}\in H^0(\widetilde{X}^\circ, \tilde{f}^*\Omega_{\widetilde{Y}^\circ}^{m}\otimes \Omega_{\widetilde{X}^\circ}^{p-m}), 
	\end{align*}
	where $\widetilde{Y}^\circ:= \pi_Y^{-1}(Y^\circ)$ and $\widetilde{X}^\circ:=\tilde{f}^{-1}(\widetilde{Y}^\circ)$.   We pick any  $y_0\in  {Y}^\circ$, and choose a simply connected coordinate open subset $(V;w_1\ldots,w_m)$ centered at $y_0$.  We  abusively denote by  $V$ a connected component of $\pi_Y^{-1}(V)$. Then $dw_1\wedge \cdots\wedge dw_m$ is a nonwhere-vanishing holomorphic $m$-form on $V$. Let $\Omega$ be a connected component of $\tilde{f}^{-1}(V)$.  We denote by $g:\Omega\to V$ the restriction of $\tilde{f}$  on $\Omega$. Then $g$ is a submersion with connected fibers. For any $y\in V$, by Step 3 in the proof of \cref{thm:vanishing2},  $dw_1\wedge \cdots\wedge dw_m$  induces  a unique holomorphic $(p-m)$-form $\alpha_{y}$ on  $g^{-1}(y)$ defined in \eqref{eq:express}  such that  $$\alpha_y\wedge \tilde{f}^*(dw_1\wedge \cdots\wedge dw_m)=\alpha.$$%within suitable coordinate open subset  $(U;z_1,\ldots,z_n)$ of $\Omega$. 
	Then by \eqref{eq:express2} and \eqref{eq:express}   together with the Fubini theorem,  we have
	\begin{align*}
		0\leq &\int_{V} \left(\int_{g^{-1}(y)}\sqrt{-1}^{(p-m)^2}\alpha_{y}\wedge \overline{\alpha_{y}}\wedge (\omega_X|_{g^{-1}(y)})^{n-p}\right) idw_1\wedge d\bar{w}_1\wedge\cdots\wedge  idw_m\wedge d\bar{w}_m\\
		&	 = \int_\Omega \sqrt{-1}^{p^2}\alpha\wedge \bar{\alpha}\wedge \omega_X^{n-p} \leq \int_{\widetilde{X}}|\alpha|^2\dvol<+\infty. 
	\end{align*} 
	Therefore,  there is a subset $Z$ of $V$ with zero Lebesgue measure,  such that for any $y\in V\backslash Z$, we have
	\begin{align}\label{eq:L2}
		\int_{g^{-1}(y)}\sqrt{-1}^{(p-m)^2}\alpha_{y}\wedge \overline{\alpha_{y}}\wedge (\omega_X|_{g^{-1}(y)})^{n-p}<\infty.
	\end{align} 
	Thus, we   construct an $L^2$ holomorphic $(p-m)$-form $\alpha_y$ on $g^{-1}(y)$ for any $y\in V\backslash Z$, which is also $d$-closed by \cref{item:d-closed}. We note that  for any $x\in \tilde{f}^{-1}(y)$,   $\alpha_y(x)=0$ if and only if $\alpha(x)=0$.
	
	Denote \( X_y := f^{-1}(y) \). By Step 3, if we choose a general point \( y \in V \setminus Z \), then there exists a morphism \( a: X_y \to A \) to an abelian variety \( A \) such that \( \dim X_y = \dim a(X_y) \). Let \( \widetilde{X_y}' \) be a connected component of the fiber product \( X_y \times_{A} \widetilde{A} \),  where $\widetilde{A}$ is the universal cover of $A$, and let \( \pi_y': \widetilde{X_y}' \to X_y \) denote the covering map. By \cite[Lemma 5.1]{EKPR12}, the map \( g^{-1}(y) \to X_y \) is an infinite Galois covering that factors through \( \pi_y': \widetilde{X_y}' \to X_y \).
	
	The conditions of \cref{prop:abelian} are therefore satisfied, and we deduce that \( \alpha_y \equiv 0 \). Since \( y \) is a general point in \( V \setminus Z \), it follows that \( \alpha(x) = 0 \) for almost all \( x \in \Omega \). By continuity, we conclude that \( \alpha(x) = 0 \) everywhere. This yields a contradiction. Therefore, we conclude that
	\[
	H^{p,0}_{(2)}(\widetilde{X}) = 0 \quad \text{for all } p \in \{0, \ldots, n-1\}.
	\]
	The theorem is proved.
\end{proof} 
We will apply \cref{thm:main2} to prove \cref{main:kollar}.
\begin{proof} [Proof of \cref{main:kollar}]
	We fix a smooth K\"ahler metric $\omega$ on $X$. 	Let 	$\cH_{(2)}^{n,q}(\widetilde{X})	$ be the $L^2$-harmonic $(n,q)$-forms with respect to the metric $\pi_X^*\omega$. By the Lefschetz theorem in  \cite[Theorem 1.2.A]{Gr}),  for any $q\in \{1,\ldots,n\}$, 
	\begin{align*}
		\cH_{(2)}^{n-q,0}(\widetilde{X})&\to  	\cH_{(2)}^{n,q}(\widetilde{X})\\\alpha&\mapsto \omega^q\wedge\alpha
	\end{align*}
	is an isomorphism. Since we have an isomorphism $\cH_{(2)}^{n,q}(\widetilde{X})\simeq H_{(2)}^{n,q}(\widetilde{X})$,  this establishes that $ H_{(2)}^{n,q}(\widetilde{X})=0$ for $q\in \{1,\ldots,n\}$. 
	
	
	We denote by $\Gamma=\pi_1(X)$ and $\dim_{\Gamma}H^{n,q}_{(2)}(\widetilde{X})$ the Von Neumann dimension of $H^{n,q}_{(2)}(\widetilde{X})$ (cf. \cite{Ati76} for the definition). 
	By Atiyah's $L^2$-index theorem along with \cref{item:vanishing}, we have
	\begin{align} \label{eq:euler}	\chi(X,K_X)=\sum_{q=0}^{n}(-1)^q\dim_{\Gamma}H^{n,q}_{(2)}(\widetilde{X})= \dim_{\Gamma}H^{n,0}_{(2)}(\widetilde{X})\geq 0.
	\end{align} 
	\Cref{item:Euler} is proved. 
	
	
	If the strict  inequality   \eqref{eq:euler}  holds, then $H^{n,0}_{(2)}(\widetilde{X})\neq 0$. We then apply \cite[Corollary 13.10]{Kol95} to conclude that $K_X$ is big. \Cref{item:strict} follows.  The theorem is proved. 
\end{proof}  
 
   
   %\bibliographystyle{ssmfalpha}
  % \bibliography{biblio2}

\begin{thebibliography}{BDDM22}
	\newcommand{\enquote}[1]{``#1''}
	\providecommand{\url}[1]{\texttt{#1}}
	\providecommand{\urlprefix}{URL }
	\providecommand{\bibinfo}[2]{#2}
	\providecommand{\eprint}[2][]{\url{#2}}
	
	\bibitem[Ati76]{Ati76}
	\bibinfo{author}{M.~F. Atiyah}.
	\newblock \enquote{\bibinfo{title}{Elliptic operators, discrete groups and von
			{N}eumann algebras}}.
	\newblock \enquote{\bibinfo{booktitle}{Colloque ``{A}nalyse et {T}opologie'' en
			l'{H}onneur de {H}enri {C}artan ({O}rsay, 1974)}}, vol. \bibinfo{volume}{No.
		32-33} of \emph{\bibinfo{series}{Ast\'{e}risque}}, \bibinfo{pages}{43--72}.
	\bibinfo{publisher}{Soc. Math. France, Paris}
	(\bibinfo{year}{1976}):\hspace{0pt}.
	
	\bibitem[BCDT24]{BCDT24}
	\bibinfo{author}{F.~{Bei}}, \bibinfo{author}{B.~{Claudon}},
	\bibinfo{author}{S.~{Diverio}} \& \bibinfo{author}{S.~{Trapani}}.
	\newblock \enquote{\bibinfo{title}{{Weakly K{\"a}hler hyperbolicity is
				birational}}}.
	\newblock \emph{\bibinfo{journal}{arXiv e-prints}},
	\textbf{(\bibinfo{year}{2024})}:\bibinfo{eid}{arXiv:2406.01734}.
	\newblock \eprint{2406.01734},
	\urlprefix\url{http://dx.doi.org/10.48550/arXiv.2406.01734}.
	
	\bibitem[BDDM22]{BDDM}
	\bibinfo{author}{D.~{Brotbek}}, \bibinfo{author}{G.~{Daskalopoulos}},
	\bibinfo{author}{Y.~{Deng}} \& \bibinfo{author}{C.~{Mese}}.
	\newblock \enquote{\bibinfo{title}{{Pluriharmonic maps into buildings and
				symmetric differentials}}}.
	\newblock \emph{\bibinfo{journal}{arXiv e-prints}},
	\textbf{(\bibinfo{year}{2022})}:\bibinfo{eid}{arXiv:2206.11835}.
	\newblock \eprint{2206.11835},
	\urlprefix\url{http://dx.doi.org/10.48550/arXiv.2206.11835}.
	
	\bibitem[BDET24]{BDET22}
	\bibinfo{author}{F.~Bei}, \bibinfo{author}{S.~Diverio},
	\bibinfo{author}{P.~Eyssidieux} \& \bibinfo{author}{S.~Trapani}.
	\newblock \enquote{\bibinfo{title}{Weakly {K}\"{a}hler hyperbolic manifolds and
			the {G}reen-{G}riffiths-{L}ang conjecture}}.
	\newblock \emph{\bibinfo{journal}{J. Reine Angew. Math.}},
	\textbf{\bibinfo{volume}{807}(\bibinfo{year}{2024})}:\bibinfo{pages}{257--297}.
	\newblock \urlprefix\url{http://dx.doi.org/10.1515/crelle-2023-0094}.
	
	\bibitem[BDIP02]{DIP02}
	\bibinfo{author}{J.~Bertin}, \bibinfo{author}{J.-P. Demailly},
	\bibinfo{author}{L.~Illusie} \& \bibinfo{author}{C.~Peters}.
	\newblock \emph{\bibinfo{title}{Introduction to {H}odge theory}},
	vol.~\bibinfo{volume}{8} of \emph{\bibinfo{series}{SMF/AMS Texts and
			Monographs}}.
	\newblock \bibinfo{publisher}{American Mathematical Society, Providence, RI;
		Soci\'{e}t\'{e} Math\'{e}matique de France, Paris} (\bibinfo{year}{2002}).
	\newblock \bibinfo{note}{Translated from the 1996 French original by James
		Lewis and Peters}.
	
	\bibitem[Car87]{Car87}
	\bibinfo{author}{J.~A. Carlson}.
	\newblock \enquote{\bibinfo{title}{The geometry of the extension class of a
			mixed {H}odge structure}}.
	\newblock \enquote{\bibinfo{booktitle}{Algebraic geometry, {B}owdoin, 1985
			({B}runswick, {M}aine, 1985)}}, vol. \bibinfo{volume}{46, Part 2} of
	\emph{\bibinfo{series}{Proc. Sympos. Pure Math.}}, \bibinfo{pages}{199--222}.
	\bibinfo{publisher}{Amer. Math. Soc., Providence, RI}
	(\bibinfo{year}{1987}):\hspace{0pt}\urlprefix\url{http://dx.doi.org/10.1090/pspum/046.2/927981}.
	
	\bibitem[CD01]{CD01}
	\bibinfo{author}{F.~Campana} \& \bibinfo{author}{J.-P. Demailly}.
	\newblock \enquote{\bibinfo{title}{Cohomologie {$L^2$} sur les rev\^etements
			d'une vari\'et\'e{} complexe compacte}}.
	\newblock \emph{\bibinfo{journal}{Ark. Mat.}},
	\textbf{\bibinfo{volume}{39}(\bibinfo{year}{2001})(\bibinfo{number}{2})}:\bibinfo{pages}{263--282}.
	\newblock \urlprefix\url{http://dx.doi.org/10.1007/BF02384557}.
	
	\bibitem[CDY22]{CDY22}
	\bibinfo{author}{B.~{Cadorel}}, \bibinfo{author}{Y.~{Deng}} \&
	\bibinfo{author}{K.~{Yamanoi}}.
	\newblock \enquote{\bibinfo{title}{{Hyperbolicity and fundamental groups of
				complex quasi-projective varieties}}}.
	\newblock \emph{\bibinfo{journal}{arXiv e-prints}},
	\textbf{(\bibinfo{year}{2022})}:\bibinfo{eid}{arXiv:2212.12225}.
	\newblock \eprint{2212.12225},
	\urlprefix\url{http://dx.doi.org/10.48550/arXiv.2212.12225}.
	
	\bibitem[CX01]{CX}
	\bibinfo{author}{J.~Cao} \& \bibinfo{author}{F.~Xavier}.
	\newblock \enquote{\bibinfo{title}{K\"{a}hler parabolicity and the {E}uler
			number of compact manifolds of non-positive sectional curvature}}.
	\newblock \emph{\bibinfo{journal}{Math. Ann.}},
	\textbf{\bibinfo{volume}{319}(\bibinfo{year}{2001})(\bibinfo{number}{3})}:\bibinfo{pages}{483--491}.
	\newblock \urlprefix\url{http://dx.doi.org/10.1007/PL00004444}.
	
	\bibitem[Dem12]{Dembook}
	\bibinfo{author}{J.-P. Demailly}.
	\newblock \emph{\bibinfo{title}{Complex Analytic and Differential Geometry}}.
	\newblock \bibinfo{publisher}{online book} (\bibinfo{year}{2012}).
	\newblock
	\urlprefix\url{http://www-fourier.ujf-grenoble.fr/~demailly/manuscripts/agbook.pdf}.
	
	\bibitem[DM24]{DM24}
	\bibinfo{author}{Y.~{Deng}} \& \bibinfo{author}{C.~{Mese}}.
	\newblock \enquote{\bibinfo{title}{{Existence and unicity of pluriharmonic maps
				to Euclidean buildings and applications}}}.
	\newblock \emph{\bibinfo{journal}{arXiv e-prints}},
	\textbf{(\bibinfo{year}{2024})}:\bibinfo{eid}{arXiv:2410.07871}.
	\newblock \eprint{2410.07871},
	\urlprefix\url{http://dx.doi.org/10.48550/arXiv.2410.07871}.
	
	\bibitem[DY24]{DY24}
	\bibinfo{author}{Y.~{Deng}} \& \bibinfo{author}{K.~{Yamanoi}}.
	\newblock \enquote{\bibinfo{title}{{Linear Shafarevich Conjecture in positive
				characteristic, Hyperbolicity and Applications}}}.
	\newblock \emph{\bibinfo{journal}{arXiv e-prints}},
	\textbf{(\bibinfo{year}{2024})}:\bibinfo{eid}{arXiv:2403.16199}.
	\newblock \eprint{2403.16199},
	\urlprefix\url{http://dx.doi.org/10.48550/arXiv.2403.16199}.
	
	\bibitem[DYK23]{DY23}
	\bibinfo{author}{Y.~{Deng}}, \bibinfo{author}{K.~{Yamanoi}} \&
	\bibinfo{author}{L.~{Katzarkov}}.
	\newblock \enquote{\bibinfo{title}{{Reductive Shafarevich Conjecture}}}.
	\newblock \emph{\bibinfo{journal}{arXiv e-prints}},
	\textbf{(\bibinfo{year}{2023})}:\bibinfo{eid}{arXiv:2306.03070}.
	\newblock \eprint{2306.03070},
	\urlprefix\url{http://dx.doi.org/10.48550/arXiv.2306.03070}.
	
	\bibitem[EKPR12]{EKPR12}
	\bibinfo{author}{P.~Eyssidieux}, \bibinfo{author}{L.~Katzarkov},
	\bibinfo{author}{T.~Pantev} \& \bibinfo{author}{M.~Ramachandran}.
	\newblock \enquote{\bibinfo{title}{Linear {S}hafarevich conjecture}}.
	\newblock \emph{\bibinfo{journal}{Ann. of Math. (2)}},
	\textbf{\bibinfo{volume}{176}(\bibinfo{year}{2012})(\bibinfo{number}{3})}:\bibinfo{pages}{1545--1581}.
	\newblock \urlprefix\url{http://dx.doi.org/10.4007/annals.2012.176.3.4}.
	
	\bibitem[Eys97]{Esy97}
	\bibinfo{author}{P.~Eyssidieux}.
	\newblock \enquote{\bibinfo{title}{La caract\'{e}ristique d'{E}uler du complexe
			de {G}auss-{M}anin}}.
	\newblock \emph{\bibinfo{journal}{J. Reine Angew. Math.}},
	\textbf{\bibinfo{volume}{490}(\bibinfo{year}{1997})}:\bibinfo{pages}{155--212}.
	\newblock \urlprefix\url{http://dx.doi.org/10.1515/crll.1997.490.155}.
	
	\bibitem[Eys99]{Eys99}
	---{}---{}---.
	\newblock \enquote{\bibinfo{title}{Syst\`emes lin\'eaires adjoints {$L^2$}}}.
	\newblock \emph{\bibinfo{journal}{Ann. Inst. Fourier (Grenoble)}},
	\textbf{\bibinfo{volume}{49}(\bibinfo{year}{1999})(\bibinfo{number}{1})}:\bibinfo{pages}{vi,
		ix--x, 141--176}.
	\newblock \urlprefix\url{http://dx.doi.org/10.5802/aif.1670}.
	
	\bibitem[Eys00]{Eys00}
	---{}---{}---.
	\newblock \enquote{\bibinfo{title}{Invariants de von {N}eumann des faisceaux
			analytiques coh\'erents}}.
	\newblock \emph{\bibinfo{journal}{Math. Ann.}},
	\textbf{\bibinfo{volume}{317}(\bibinfo{year}{2000})(\bibinfo{number}{3})}:\bibinfo{pages}{527--566}.
	\newblock \urlprefix\url{http://dx.doi.org/10.1007/PL00004413}.
	
	\bibitem[{Eys}04]{Eys04}
	\bibinfo{author}{P.~{Eyssidieux}}.
	\newblock \enquote{\bibinfo{title}{{Sur la convexit\'e holomorphe des
				rev\^etements lin\'eaires r\'eductifs d'une vari\'et\'e projective
				alg\'ebrique complexe}}}.
	\newblock \emph{\bibinfo{journal}{{Invent. Math.}}},
	\textbf{\bibinfo{volume}{156}(\bibinfo{year}{2004})(\bibinfo{number}{3})}:\bibinfo{pages}{503--564}.
	\newblock \urlprefix\url{http://dx.doi.org/10.1007/s00222-003-0345-0}.
	
	\bibitem[GL87]{GL87}
	\bibinfo{author}{M.~Green} \& \bibinfo{author}{R.~Lazarsfeld}.
	\newblock \enquote{\bibinfo{title}{Deformation theory, generic vanishing
			theorems, and some conjectures of {E}nriques, {C}atanese and {B}eauville}}.
	\newblock \emph{\bibinfo{journal}{Invent. Math.}},
	\textbf{\bibinfo{volume}{90}(\bibinfo{year}{1987})(\bibinfo{number}{2})}:\bibinfo{pages}{389--407}.
	\newblock \urlprefix\url{http://dx.doi.org/10.1007/BF01388711}.
	
	\bibitem[Gro91]{Gr}
	\bibinfo{author}{M.~Gromov}.
	\newblock \enquote{\bibinfo{title}{K\"{a}hler hyperbolicity and {$L_2$}-{H}odge
			theory}}.
	\newblock \emph{\bibinfo{journal}{J. Differential Geom.}},
	\textbf{\bibinfo{volume}{33}(\bibinfo{year}{1991})(\bibinfo{number}{1})}:\bibinfo{pages}{263--292}.
	\newblock \urlprefix\url{http://projecteuclid.org/euclid.jdg/1214446039}.
	
	\bibitem[GS92]{GS92}
	\bibinfo{author}{M.~{Gromov}} \& \bibinfo{author}{R.~{Schoen}}.
	\newblock \enquote{\bibinfo{title}{{Harmonic maps into singular spaces and
				$p$-adic superrigidity for lattices in groups of rank one}}}.
	\newblock \emph{\bibinfo{journal}{{Publ. Math., Inst. Hautes \'Etud. Sci.}}},
	\textbf{\bibinfo{volume}{76}(\bibinfo{year}{1992})}:\bibinfo{pages}{165--246}.
	\newblock \urlprefix\url{http://dx.doi.org/10.1007/BF02699433}.
	
	\bibitem[GZ17]{GZ17}
	\bibinfo{author}{V.~Guedj} \& \bibinfo{author}{A.~Zeriahi}.
	\newblock \emph{\bibinfo{title}{Degenerate complex {M}onge-{A}mp\`ere
			equations}}, vol.~\bibinfo{volume}{26} of \emph{\bibinfo{series}{EMS Tracts
			in Mathematics}}.
	\newblock \bibinfo{publisher}{European Mathematical Society (EMS), Z\"urich}
	(\bibinfo{year}{2017}).
	\newblock \urlprefix\url{http://dx.doi.org/10.4171/167}.
	
	\bibitem[Hua19]{Hua19}
	\bibinfo{author}{S.~Huang}.
	\newblock \enquote{\bibinfo{title}{A note on existence of exhaustion functions
			and its applications}}.
	\newblock \emph{\bibinfo{journal}{J. Geom. Anal.}},
	\textbf{\bibinfo{volume}{29}(\bibinfo{year}{2019})(\bibinfo{number}{2})}:\bibinfo{pages}{1649--1659}.
	\newblock \urlprefix\url{http://dx.doi.org/10.1007/s12220-018-0055-x}.
	
	\bibitem[JZ00]{JZ}
	\bibinfo{author}{J.~Jost} \& \bibinfo{author}{K.~Zuo}.
	\newblock \enquote{\bibinfo{title}{Vanishing theorems for {$L^2$}-cohomology on
			infinite coverings of compact {K}\"{a}hler manifolds and applications in
			algebraic geometry}}.
	\newblock \emph{\bibinfo{journal}{Comm. Anal. Geom.}},
	\textbf{\bibinfo{volume}{8}(\bibinfo{year}{2000})(\bibinfo{number}{1})}:\bibinfo{pages}{1--30}.
	\newblock \urlprefix\url{http://dx.doi.org/10.4310/CAG.2000.v8.n1.a1}.
	
	\bibitem[Kol95]{Kol95}
	\bibinfo{author}{J.~Koll{\'a}r}.
	\newblock \emph{\bibinfo{title}{Shafarevich maps and automorphic forms}}.
	\newblock \bibinfo{publisher}{Princeton University Press},
	\bibinfo{address}{Princeton (N.J.)} (\bibinfo{year}{1995}).
	
	\bibitem[KP23]{KP23}
	\bibinfo{author}{T.~Kaletha} \& \bibinfo{author}{G.~Prasad}.
	\newblock \emph{\bibinfo{title}{Bruhat-{Tits} theory. {A} new approach (to
			appear)}}, vol.~\bibinfo{volume}{44} of \emph{\bibinfo{series}{New Math.
			Monogr.}}
	\newblock \bibinfo{publisher}{Cambridge: Cambridge University Press}
	(\bibinfo{year}{2023}).
	
	\bibitem[Rou23]{Rou23}
	\bibinfo{author}{G.~Rousseau}.
	\newblock \emph{\bibinfo{title}{Euclidean buildings---geometry and group
			actions}}, vol.~\bibinfo{volume}{35} of \emph{\bibinfo{series}{EMS Tracts in
			Mathematics}}.
	\newblock \bibinfo{publisher}{EMS Press, Berlin} (\bibinfo{year}{[2023]
		\copyright 2023}).
	\newblock \urlprefix\url{http://dx.doi.org/10.4171/etm/35}.
	
	\bibitem[Shi72]{Shi72}
	\bibinfo{author}{B.~Shiffman}.
	\newblock \enquote{\bibinfo{title}{Extension of positive line bundles and
			meromorphic maps}}.
	\newblock \emph{\bibinfo{journal}{Invent. Math.}},
	\textbf{\bibinfo{volume}{15}(\bibinfo{year}{1972})(\bibinfo{number}{4})}:\bibinfo{pages}{332--347}.
	\newblock \urlprefix\url{http://dx.doi.org/10.1007/BF01405087}.
	
	\bibitem[Tak04]{Tak04}
	\bibinfo{author}{S.~Takayama}.
	\newblock \enquote{\bibinfo{title}{On the existence of pluricanonical forms on
			varieties with infinite fundamental group}}.
	\newblock \emph{\bibinfo{journal}{Amer. J. Math.}},
	\textbf{\bibinfo{volume}{126}(\bibinfo{year}{2004})(\bibinfo{number}{6})}:\bibinfo{pages}{1221--1235}.
	\newblock
	\urlprefix\url{http://muse.jhu.edu/journals/american_journal_of_mathematics/v126/126.6takayama.pdf}.
	
\end{thebibliography}


\end{document}



